\documentclass{article}
\usepackage{fuzz}
\makeatletter
\newenvironment{theorem}[1]{\medskip\noindent\textbf{Theorem (#1).}\@nameuse{zed}}{\@nameuse{endzed}}
\makeatother

\begin{document}

\title{A Formal Specification of \texttt{points} (target routing state and Slice 1 refinement)}
\author{Punt Labs}
\maketitle

\section{Purpose and scope}

This is the \emph{target} specification --- the robust design a from-scratch
re-implementation targets --- paired with
\texttt{model-provider.current.spec.tex}, which models the predecessor.

The organising idea, reached by separating two axes the combined tool conflated:
\emph{routing} (how a request reaches an upstream --- base URL, auth, which
source) is session-global per harness; \emph{model-choice} (which model serves a
request) is per-agent, user-authored, and out of scope. \texttt{points}
owns routing. Its whole job is: the operator picks, per harness, a \emph{source}
to route through and a default model, and the tool writes exactly that route so
the chosen model --- and any model any agent later names --- reaches the upstream
the operator picked.

\textbf{The user picks the route.} There is no tool-defined precedence and no
ban on overlap: if two sources both serve a provider, the operator simply routes
a given harness through whichever source they prefer. \emph{That choice is the
	entire tool.} The managed state is therefore a per-harness map from harness to
source, set by the operator, never derived.

\textbf{Multiple sources.} One operator may have a single source (a gateway like
Bifrost that serves every provider and carries both protocols); another may have
several narrow sources (direct OpenAI, OpenRouter) and no gateway at all. The
spec models a \emph{set} of configured sources and a per-harness route, so the
single-gateway case is just the degenerate ``every harness routed through one
source that serves everything.''

\section{Given sets and enumerations}

A model id is an opaque catalog token (e.g.
\texttt{openrouter/anthropic/claude-opus-4.8}); its leading segment names a
provider. A provider is the upstream a source routes to (\texttt{anthropic},
\texttt{openai}, \texttt{bedrock}, \texttt{openrouter} as an aggregator). A
source is a configured route target: a gateway (Bifrost, Envoy, AgentGateway) or
a direct vendor/aggregator endpoint.

% POINTS-CHECK BEGIN types.routing-carriers type
\begin{zed}
	[ModelId, Provider, Source]
\end{zed}
% POINTS-CHECK END types.routing-carriers

% POINTS-CHECK BEGIN types.transaction-carriers type
\begin{zed}
	[JSONPointer, Path, Identity, Owner, Mode, ByteString, Delta, FileImage]
\end{zed}
% POINTS-CHECK END types.transaction-carriers

The four harnesses.

\begin{zed}
	Harness ::= claudecode | codex | pi | opencode
\end{zed}

The protocol each harness speaks. Claude Code speaks Anthropic; the other three
speak OpenAI. A source must expose a harness's protocol to carry it.

\begin{zed}
	Protocol ::= anthropic | openai
\end{zed}

Auth modes a source may accept: a per-host virtual key / API key ($vkauth$), or
the operator's Claude subscription OAuth ($subauth$).

\begin{zed}
	AuthMode ::= vkauth | subauth
\end{zed}

Boolean and completeness values are made finite explicitly so that the bounded
model is executable rather than relying on an implementation-language type.

% POINTS-CHECK BEGIN types.bool type
\begin{zed}
	Bool ::= yes | no \\
	ContentState ::= partial | complete
\end{zed}
% POINTS-CHECK END types.bool

Claude Code's built-in model-alias tiers --- the Agent tool's
\texttt{opus|sonnet|haiku|fable} knobs. Not user agent definitions; built-in
alias slots the tool sets so a model-less spawn resolves to a routable id. A
fixed, finite set.

\begin{zed}
	Tier ::= opus | sonnet | haiku | fable
\end{zed}

\section{Harness and model properties}

Which protocol each harness speaks.

\begin{axdef}
	speaks : Harness \fun Protocol
	\where
	speaks = \{ claudecode \mapsto anthropic, \\
	\t1         codex \mapsto openai, \\
	\t1         pi \mapsto openai, \\
	\t1         opencode \mapsto openai \}
\end{axdef}

Which harnesses canonicalize the configured id (strip the provider prefix),
losing the routing information unless an override restores the verbatim id. Only
Claude Code.

\begin{axdef}
	canonicalizes : \power Harness
	\where
	canonicalizes = \{ claudecode \}
\end{axdef}

Which harnesses present a native model picker that validates against a
tool-written list (so the tool must populate it). pi and opencode do; codex and
Claude Code do not.

\begin{axdef}
	showsPicker : \power Harness
	\where
	showsPicker = \{ pi, opencode \}
\end{axdef}

Every model id carries a provider --- the only thing routing depends on, since a
source fans out by provider. Total, so there are no partial-application gaps.

\begin{axdef}
	providerOf : ModelId \fun Provider
\end{axdef}

\section{The source contract}

For each source: $srcServes$ is the set of providers this account is entitled to
reach through it (a gateway serves many; a direct source one or few; the Bedrock
Anthropic-models entitlement gap is one such restriction). $srcProtocols$ is the
protocols it exposes; $srcAuth$ the auth modes it accepts. Every source serves at
least one provider, exposes at least one protocol, and accepts at least one auth
mode.

\begin{axdef}
	srcServes : Source \fun \power Provider \\
	srcProtocols : Source \fun \power Protocol \\
	srcAuth : Source \fun \power AuthMode
	\where
	\forall s : Source @ \\
	\t1 srcServes~s \neq \emptyset \land \\
	\t1 srcProtocols~s \neq \emptyset \land \\
	\t1 srcAuth~s \neq \emptyset
\end{axdef}

The universe contains at least one \emph{Anthropic-capable} source with a
routable model and both auth modes --- a witness, not a requirement on any
particular configuration. Its sole purpose is to force the model checker's
constant search to admit a source through which Claude Code can route, so the
Claude-Code guarantees (override, four-alias routability, the $subauth$ path) are
exercised on reachable states rather than passing vacuously. Operators with no
Anthropic source simply never route Claude Code; this witness does not oblige
them to.

\begin{axdef}
	anthropicWitness : Source \\
	anthropicModel : ModelId
	\where
	anthropic \in srcProtocols~anthropicWitness \\
	openai \in srcProtocols~anthropicWitness \\
	srcAuth~anthropicWitness = \{ vkauth, subauth \} \\
	providerOf~anthropicModel \in srcServes~anthropicWitness
\end{axdef}

The picker pattern: the set of ids the operator's picker expression matches (a
regex in the implementation; its match set here). The native picker offers the
pattern intersected with what is routable on that harness's source --- a view,
never a constraint on what may be routed.

\begin{axdef}
	matchesPattern : \power ModelId
\end{axdef}

The selectable catalog is stricter than provider routability.  This distinction
prevents a same-provider but unsupported identifier from entering managed
configuration merely because its provider is served.

\begin{axdef}
	selectableModels : Source \fun \power ModelId
	\where
	\forall s : Source @ \forall m : selectableModels~s @
	providerOf~m \in srcServes~s
\end{axdef}

\section{Raw, coherent, and observed configuration}

The raw managed projection is routing, per harness, plus the fixtures the one
canonicalizing harness needs. $configured$ is the set of sources the operator
has set up. $routeVia$ maps each \emph{routed} harness to the source the operator
chose for it; its domain is exactly the harnesses currently on a managed route
(others are on their vendor default / direct). $model$ gives each routed
harness its default model; $hauth$ its auth mode. $override$ records which
canonicalizing harnesses have the \texttt{modelOverrides} fixture written.
$aliases$ gives Claude Code's four tier knobs. $shown$ gives, for each
picker-bearing routed harness, the ids its native picker offers.

Agent definitions are deliberately absent: the tool does not model, read, or
write them. A per-agent model pin is orthogonal user state that routes through
the same per-harness route.

The invariant folds in \emph{all} route coherence, so every reachable state
--- not merely states an operation happened to build --- is checked by the model
checker:

Unlike the former model, raw state carries no coherence invariant.  It can
therefore represent malformed and partially managed external observations and
make verification failures reachable.  $registered$ is only the certified
managed registration projection; it is not every native declaration preserved
in an OpenCode document.  Registration and the optional picker view are
different concerns.

\begin{schema}{RawConfig}
	configured : \power Source \\
	routeVia : Harness \pfun Source \\
	model : Harness \pfun ModelId \\
	hauth : Harness \pfun AuthMode \\
	override : \power Harness \\
	aliases : Tier \pfun ModelId \\
	registered : Harness \pfun \power ModelId \\
	shown : Harness \pfun \power ModelId
\end{schema}

\begin{axdef}
	needsRegistration : \power Harness
	\where
	needsRegistration = \{ pi, opencode \}
\end{axdef}

The coherent subtype is the only place that assumes routing invariants.  In
particular subscription auth belongs only to Claude Code, not merely to any
source which happens to expose the Anthropic protocol.

\begin{schema}{Config}
	RawConfig
	\where
	\dom model = \dom routeVia \\
	\dom hauth = \dom routeVia \\
	\ran routeVia \subseteq configured \\
	override \subseteq canonicalizes \cap \dom routeVia \\
	\dom registered = needsRegistration \cap \dom routeVia \\
	\dom shown = showsPicker \cap \dom routeVia \\
	( \forall h : \dom routeVia @ \\
	\t1 speaks~h \in srcProtocols~(routeVia~h) \land \\
	\t1 providerOf~(model~h) \in srcServes~(routeVia~h) \land \\
	\t1 model~h \in selectableModels~(routeVia~h) \land \\
	\t1 hauth~h \in srcAuth~(routeVia~h) \land \\
	\t1 ( h \in canonicalizes \implies h \in override ) \land \\
	\t1 ( hauth~h = subauth \implies h = claudecode ) ) \\
	( claudecode \in \dom routeVia \implies \\
	\t1 ( \dom aliases = \{ opus, sonnet, haiku, fable \} \land \\
	\t1   ( \forall t : \dom aliases @ \\
	\t2     providerOf~(aliases~t) \in srcServes~(routeVia~claudecode) ) ) ) \\
	( claudecode \notin \dom routeVia \implies aliases = \emptyset ) \\
	( \forall h : \dom shown @ \\
	\t1 ( \forall m : shown~h @ \\
	\t2 m \in matchesPattern \land m \in selectableModels~(routeVia~h) ) ) \\
	( \forall h : \dom registered @ \\
	\t1 model~h \in registered~h \land \\
	\t1 registered~h \subseteq selectableModels~(routeVia~h) )
\end{schema}

External observation is deliberately not a $Config$.  The decoder is total and
classifies absent, unreadable, malformed, decoded, missing-slot, invalid-slot,
and present-slot inputs.  The raw fields remain available only as decoded data;
every predicate below guards partial-function application with a domain test.

\begin{zed}
	ObservationKind ::= obsAbsent | obsUnreadable | obsMalformed | obsDecoded \\
	SlotState ::= slotMissing | slotInvalid | slotPresent
\end{zed}

\begin{schema}{ObservedConfig}
	RawConfig \\
	observationKind : ObservationKind \\
	routeSlot, modelSlot, authSlot, registrationSlot : SlotState
	\where
	observationKind \neq obsDecoded \implies \\
	\t1 routeSlot = slotMissing \land modelSlot = slotMissing \land \\
	\t1 authSlot = slotMissing \land registrationSlot = slotMissing
\end{schema}

\section{Initial state}

Never configured: no routed harness, no override, no aliases, no picker lists.

\begin{schema}{Init}
	RawConfig'
	\where
	routeVia' = \emptyset \\
	model' = \emptyset \\
	hauth' = \emptyset \\
	override' = \emptyset \\
	aliases' = \emptyset \\
	registered' = \emptyset \\
	shown' = \emptyset
\end{schema}

\section{The core guarantees}

Because coherence lives in $Config$, these schemas \emph{name} the guarantees
for the reader and for targeted checking; they are not the only thing enforcing
them. Each is falsifiable: a state violating it is not a legal $Config$.

\textbf{Default routability (per harness).} Every routed harness's default model
reaches its source's providers --- the operator's pick is served, per harness,
by the source the operator chose. This is about the \emph{selected} model, not a
definitional set, so it has teeth: a model whose provider the source does not
serve violates it.

\begin{schema}{DefaultRoutable}
	Config
	\where
	\forall h : \dom routeVia @ \\
	\t1 providerOf~(model~h) \in srcServes~(routeVia~h)
\end{schema}

\textbf{Protocol coherence.} Every routed harness's source speaks that harness's
protocol --- Claude Code only ever routes through an Anthropic-capable source,
the OpenAI harnesses through OpenAI-capable ones.

\begin{schema}{ProtocolCoherent}
	Config
	\where
	\forall h : \dom routeVia @ speaks~h \in srcProtocols~(routeVia~h)
\end{schema}

\textbf{Claude routing-preservation.} When Claude Code is routed, its override is
written (so its canonicalized id is restored on the wire) and all four tier
knobs are set to ids routable on its source (so no model-less tier spawn
mis-routes). The ``all four set'' requirement is what stops a routed-but-empty
alias state passing as correct.

\begin{schema}{ClaudeRoutingPreserved}
	Config
	\where
	claudecode \in \dom routeVia \implies \\
	\t1 ( claudecode \in override \land \\
	\t1   \dom aliases = \{ opus, sonnet, haiku, fable \} \land \\
	\t1   ( \forall t : \dom aliases @ \\
	\t2     providerOf~(aliases~t) \in srcServes~(routeVia~claudecode) ) )
\end{schema}

\textbf{Picker soundness.} A native picker never offers an unroutable id: every
shown id matches the pattern and its provider is served by that harness's
source. The picker is a view ($shown~h$ need not contain, nor be contained by,
any particular model --- the operator may route an id the picker does not show).
An \emph{empty} picker view ($shown~h = \emptyset$) is a legal, intentional state
(e.g. immediately after routing, before a pattern is applied); the guarantee is
``every id the picker shows is routable,'' which an empty view satisfies
correctly --- it withholds nothing unroutable because it shows nothing.

\begin{schema}{PickerSound}
	Config
	\where
	\forall h : \dom shown @ \forall m : shown~h @ \\
	\t1 m \in matchesPattern \land providerOf~m \in srcServes~(routeVia~h)
\end{schema}

\textbf{The forbidden failure.} A routed harness whose default's provider its
source does not serve, a canonicalizing harness with no override, or Claude Code
routed with a tier knob unset or unroutable --- any request that reaches no valid
upstream. Not a legal $Config$; named as the predicate the invariant excludes.

\begin{schema}{SilentMisroute}
	configured : \power Source \\
	routeVia : Harness \pfun Source \\
	model : Harness \pfun ModelId \\
	hauth : Harness \pfun AuthMode \\
	override : \power Harness \\
	aliases : Tier \pfun ModelId \\
	shown : Harness \pfun \power ModelId
	\where
	\dom model = \dom routeVia \\
	( ( \exists h : \dom routeVia @ \\
	\t1 providerOf~(model~h) \notin srcServes~(routeVia~h) ) \lor \\
	\t1 ( \exists h : canonicalizes \cap \dom routeVia @ h \notin override ) \lor \\
	\t1 ( claudecode \in \dom routeVia \land \\
	\t2 \dom aliases \neq \{ opus, sonnet, haiku, fable \} ) )
\end{schema}

\section{Factored operation: \texttt{route}}

The one selecting operation. The operator names a set of harnesses $hs?$, a
source $s?$, a default model $m?$, an auth mode $a?$, and (when Claude Code is in
$hs?$) the tier ids $tiers?$. It routes exactly those harnesses through $s?$,
last-write-wins per harness --- so a later \texttt{route} of a different source
for a harness is the operator changing their preferred route, with no precedence
rule and no overlap ban. Preconditions (all enforced \emph{in the operation}, so
the shared effect alone establishes $Config'$):
$s?$ is configured; $s?$ speaks every chosen harness's protocol; $m?$'s provider
is served by $s?$; $a?$ is accepted by $s?$; $subauth$ only where $s?$ is
Anthropic-capable; and if Claude Code is chosen, all four tiers are supplied and
routable on $s?$.

The operation is factored so its effect cannot smuggle $Config'$ in as a
hypothesis.  Empty selection, incoherent pre-state, and every eligibility rule
are explicit preconditions.  The effect assigns every raw field.

\begin{schema}{RouteInput}
	hs? : \power Harness \\
	s? : Source \\
	m? : ModelId \\
	a? : AuthMode \\
	tiers? : Tier \pfun ModelId
\end{schema}

\begin{schema}{RoutePre}
	Config \\
	RouteInput
	\where
	hs? \neq \emptyset \\
	s? \in configured \\
	( \forall h : hs? @ speaks~h \in srcProtocols~s? ) \\
	m? \in selectableModels~s? \\
	a? \in srcAuth~s? \\
	( a? = subauth \implies hs? = \{ claudecode \} ) \\
	( claudecode \in hs? \implies \\
	\t1 ( \dom tiers? = \{ opus, sonnet, haiku, fable \} \land \\
	\t1   ( \forall t : \dom tiers? @ tiers?~t \in selectableModels~s? ) ) )
\end{schema}

\begin{schema}{RouteEffect}
	\Delta RawConfig \\
	RouteInput
	\where
	configured' = configured \\
	routeVia' = routeVia \oplus \{ h : hs? @ h \mapsto s? \} \\
	model' = model \oplus \{ h : hs? @ h \mapsto m? \} \\
	hauth' = hauth \oplus \{ h : hs? @ h \mapsto a? \} \\
	override' = override \cup ( canonicalizes \cap hs? ) \\
	( claudecode \in hs? \implies aliases' = tiers? ) \\
	( claudecode \notin hs? \implies aliases' = aliases ) \\
	registered' = ( needsRegistration \cap \dom routeVia' ) \dres registered \\
	\t1 \oplus \{ h : needsRegistration \cap hs? @ h \mapsto \{ m? \} \} \\
	shown' = ( showsPicker \cap \dom routeVia' ) \dres shown \\
	\t1 \oplus \{ h : showsPicker \cap hs? @ h \mapsto \emptyset \}
\end{schema}

\begin{schema}{RouteOk}
	RoutePre \\
	RouteEffect
\end{schema}

The refusals, each a no-op with a distinct code: source not configured; source
cannot carry a chosen harness's protocol; the default's provider not served; the
auth mode not accepted; a $subauth$ request to a non-Anthropic source; or Claude
Code chosen without all four routable tiers.

\begin{zed}
	RouteReason ::= emptySelection | incoherentPrestate | sourceNotConfigured | \\
	\t1 protocolMismatch | modelNotSelectable | authNotAccepted | \\
	\t1 subscriptionNotClaudeOnly | tiersIncomplete | tierNotSelectable
\end{zed}

\begin{schema}{RouteRefused}
	\Xi RawConfig \\
	RouteInput \\
	reasons! : \power RouteReason
	\where
	( emptySelection \in reasons! \iff hs? = \emptyset ) \\
	( incoherentPrestate \in reasons! \iff \theta RawConfig \notin Config ) \\
	( sourceNotConfigured \in reasons! \iff s? \notin configured ) \\
	( protocolMismatch \in reasons! \iff \\
	\t1 ( \exists h : hs? @ speaks~h \notin srcProtocols~s? ) ) \\
	( modelNotSelectable \in reasons! \iff m? \notin selectableModels~s? ) \\
	( authNotAccepted \in reasons! \iff a? \notin srcAuth~s? ) \\
	( subscriptionNotClaudeOnly \in reasons! \iff \\
	\t1 a? = subauth \land hs? \neq \{ claudecode \} ) \\
	( tiersIncomplete \in reasons! \iff claudecode \in hs? \land \\
	\t1 \dom tiers? \neq \{ opus, sonnet, haiku, fable \} ) \\
	( tierNotSelectable \in reasons! \iff claudecode \in hs? \land \\
	\t1 ( \exists t : \dom tiers? @ tiers?~t \notin selectableModels~s? ) ) \\
	reasons! \neq \emptyset
\end{schema}

\begin{zed}
	Route \defs RouteOk \lor RouteRefused
\end{zed}

\section{Factored operation: \texttt{unroute} (direct)}

Return a nonempty set of harnesses to their vendor defaults.  Like route, the
general-target operation separates input, coherent precondition, complete raw
effect, and exact no-op refusal.  Malformed observations are for read commands;
they never silently authorize mutation.  Removing registration is part of the
effect.  Clearing Claude Code also clears aliases.

\begin{schema}{UnrouteInput}
	hs? : \power Harness
\end{schema}

\begin{schema}{UnroutePre}
	Config \\
	UnrouteInput
	\where
	hs? \neq \emptyset
\end{schema}

\begin{schema}{UnrouteEffect}
	\Delta RawConfig \\
	UnrouteInput
	\where
	configured' = configured \\
	routeVia' = hs? \ndres routeVia \\
	model' = hs? \ndres model \\
	hauth' = hs? \ndres hauth \\
	override' = override \setminus hs? \\
	( claudecode \in hs? \implies aliases' = \emptyset ) \\
	( claudecode \notin hs? \implies aliases' = aliases ) \\
	registered' = hs? \ndres registered \\
	shown' = hs? \ndres shown
\end{schema}

\begin{schema}{UnrouteOk}
	UnroutePre \\
	UnrouteEffect
\end{schema}

\begin{zed}
	UnrouteReason ::= unrouteEmptySelection | unrouteIncoherentPrestate
\end{zed}

\begin{schema}{UnrouteRefused}
	\Xi RawConfig \\
	UnrouteInput \\
	unrouteReasons! : \power UnrouteReason
	\where
	( unrouteEmptySelection \in unrouteReasons! \iff hs? = \emptyset ) \\
	( unrouteIncoherentPrestate \in unrouteReasons! \iff
	\theta RawConfig \notin Config ) \\
	unrouteReasons! \neq \emptyset
\end{schema}

\begin{zed}
	Unroute \defs UnrouteOk \lor UnrouteRefused
\end{zed}

\begin{theorem}{UnrouteEstablishesConfig}
	\forall \Delta RawConfig; hs? : \power Harness @ \\
	\t1 UnroutePre \land UnrouteEffect \implies \theta RawConfig' \in Config
\end{theorem}

\section{Read-only operations and non-vacuous verification}

$Status$ reports the per-harness route and default. $List$ reports, for a source,
the providers it serves. $Verify$ is the probe: $ok! = 1$ exactly when the
configuration is a legal $Config$ in which every routed harness is fully correct
--- default routable, protocol matched, auth accepted, canonicalizing override
present, and (if Claude Code routed) all four tiers routable. Because $Config$
already enforces these, $Verify$ cannot pass a mis-routed state.

\begin{schema}{Status}
	\Xi Config \\
	route! : Harness \pfun Source \\
	dflt! : Harness \pfun ModelId
	\where
	route! = routeVia \\
	dflt! = model
\end{schema}

\begin{schema}{List}
	\Xi Config \\
	s? : Source \\
	provs! : \power Provider
	\where
	provs! = srcServes~s?
\end{schema}

\begin{zed}
	ProbeOutcome ::= probeNotRun | probeCredentialMissing | probeCredentialInvalid | \\
	\t1 probeTimeout | probeUnreachable | probeRedirect | probeAuthFailure | \\
	\t1 probeHTTPFailure | probeInvalidCatalog | probeModelAbsent | probePresent
\end{zed}

\begin{schema}{VerifyInput}
	ObservedConfig \\
	probe? : ProbeOutcome
\end{schema}

\begin{schema}{LocalPass}
	ObservedConfig
	\where
	observationKind = obsDecoded \\
	routeSlot = slotPresent \\
	modelSlot = slotPresent \\
	authSlot = slotPresent \\
	registrationSlot = slotPresent \\
	\theta RawConfig \in Config \\
	\dom routeVia \neq \emptyset
\end{schema}

\begin{schema}{LivePass}
	VerifyInput
	\where
	probe? = probePresent
\end{schema}

\begin{schema}{VerifyPass}
	\Xi ObservedConfig \\
	probe? : ProbeOutcome \\
	ok! : Bool
	\where
	LocalPass \\
	probe? = probePresent \\
	ok! = yes
\end{schema}

\begin{schema}{VerifyFail}
	\Xi ObservedConfig \\
	probe? : ProbeOutcome \\
	ok! : Bool
	\where
	\lnot ( LocalPass \land probe? = probePresent ) \\
	ok! = no
\end{schema}

\begin{zed}
	Verify \defs VerifyPass \lor VerifyFail
\end{zed}

\section{The system invariant and the apply obligation}

$SystemOK$ restates, as a single predicate, that every routed harness is fully
correct --- it is exactly the body $Config$ enforces, named for checking.

\begin{schema}{SystemOK}
	Config
	\where
	( \forall h : \dom routeVia @ \\
	\t1 providerOf~(model~h) \in srcServes~(routeVia~h) \land \\
	\t1 speaks~h \in srcProtocols~(routeVia~h) \land \\
	\t1 hauth~h \in srcAuth~(routeVia~h) \land \\
	\t1 ( h \in canonicalizes \implies h \in override ) ) \land \\
	( claudecode \in \dom routeVia \implies \\
	\t1 \dom aliases = \{ opus, sonnet, haiku, fable \} )
\end{schema}

The real obligation, with no smuggled hypothesis: every harness $RouteOk$
touches ends routable on the source the operator chose, established by $RouteOk$'s
own preconditions and effect alone. ProB checks it via \texttt{-cbc\_assertions}.

\begin{theorem}{RouteEstablishesRoutability}
	\forall \Delta RawConfig; hs? : \power Harness; s? : Source; m? : ModelId; \\
	\t1 a? : AuthMode; tiers? : Tier \pfun ModelId @ \\
	\t2 RouteOk \implies \\
	\t3 ( \forall h : hs? @ \\
	\t4 providerOf~(model'~h) \in srcServes~(routeVia'~h) )
\end{theorem}

\begin{theorem}{RouteEstablishesConfig}
	\forall \Delta RawConfig; hs? : \power Harness; s? : Source; m? : ModelId; \\
	\t1 a? : AuthMode; tiers? : Tier \pfun ModelId @ \\
	\t2 RoutePre \land RouteEffect \implies \theta RawConfig' \in Config
\end{theorem}

\section{Slice 1: explicit-file OpenCode refinement}

The first delivery slice is intentionally narrower than the general model.  It
supports one harness/source pair, an explicitly named strict-JSON file, an
explicit Bifrost deployment root, and exactly two model tokens.  This is a
refinement, not a claim that the general source model has only one source.

\begin{axdef}
	sol, astra, unsupported : ModelId \\
	bifrostFixture : Source
	\where
	sol \neq astra \\
	unsupported \notin \{ sol, astra \} \\
	providerOf~unsupported = providerOf~sol \\
	selectableModels~bifrostFixture = \{ sol, astra \} \\
	openai \in srcProtocols~bifrostFixture \\
	vkauth \in srcAuth~bifrostFixture
\end{axdef}

\begin{zed}
	RootState ::= validHTTPSRoot | validLoopbackHTTPRoot | validInsecureHTTPRoot | \\
	\t1 invalidRoot | endpointFormInsteadOfRoot \\
	SliceHarness ::= sliceOpenCode | otherHarnessToken \\
	SliceSource ::= sliceBifrost | otherSourceToken
\end{zed}

\begin{schema}{ConcreteInputGate}
	sliceHarness? : SliceHarness \\
	sliceSource? : SliceSource \\
	sliceModel? : ModelId \\
	rootState? : RootState \\
	configPathLexicallyAbsolute?, insecureConsent?, gateAccepted! : Bool
	\where
	gateAccepted! = yes \iff \\
	\t1 sliceHarness? = sliceOpenCode \land sliceSource? = sliceBifrost \land \\
	\t1 sliceModel? \in \{ sol, astra \} \land \\
	\t1 configPathLexicallyAbsolute? = yes \land \\
	\t1 ( rootState? \in \{ validHTTPSRoot, validLoopbackHTTPRoot \} \lor \\
	\t2 ( rootState? = validInsecureHTTPRoot \land insecureConsent? = yes ) )
\end{schema}

The concrete spellings are respectively
\texttt{bedrock/gpt-5.6-sol} and \texttt{openai/gpt-6-astra}.  A bounded
valuation must give $unsupported$ the same provider as one supported model;
that witnesses that provider membership alone does not imply fixture
membership.  Route is offline.  Status reads only the explicit file.  Verify
checks that file and then attempts at most one authenticated catalog request.

For deployment root $R$, the managed OpenCode base is $R$/\texttt{openai/v1},
the catalog is $R$/\texttt{v1/models}, and the Responses endpoint is beneath the
managed base.  The root must be absolute HTTP(S), have a nonempty host and valid
optional port, and contain no userinfo, query, fragment, controls, backslashes,
dot segments, encoded separators, or encoded dot-segment ambiguity.  Supplying
an endpoint form in place of the root is rejected.  One trailing slash is
removed.  HTTPS and literal loopback HTTP are accepted; exact, case-normalized
\texttt{localhost} HTTP is accepted only with loopback-only connection.  Every
other HTTP use requires fresh \texttt{--allow-insecure-http} consent.  Consent
is not persisted, never weakens TLS, and never enables redirects.

\subsection{Native OpenCode V2 projection and ownership}

For selected token $M$ the native document has root default
\texttt{bifrost/M}; \texttt{providers.bifrost.package} is
\texttt{@opencode/ai/providers/openai-compatible/responses};
\texttt{env} is exactly \texttt{["BIFROST\_VK"]}; settings contain the derived
\texttt{baseURL} and literal \texttt{\{env:BIFROST\_VK\}}; and model $M$ is
registered with \texttt{modelID=M}, \texttt{disabled=false}, and tool
capability true.  Registration does not prune another registration and is not a
picker filter.  No cost, limit, variant, legacy singular \texttt{provider},
\texttt{options}, \texttt{agent}, or \texttt{small\_model} field is invented.

One pure normalizer returns a new document, sorted unique changed RFC~6901
pointers, and sorted unique conflict pointers.  It sets the owned projection;
removes provider \texttt{canonical}; removes model/variant package, endpoint,
credential, provider, route/auth-control header, and request-body model
overrides; and applies that overlay removal to every existing Bifrost model and
declared variant.  Controlled headers are case-insensitive
\texttt{authorization}, \texttt{proxy-authorization}, \texttt{host},
\texttt{api-key}, \texttt{x-api-key}, \texttt{openai-organization},
\texttt{openai-project}, \texttt{x-provider-api-keys}, and prefixes
\texttt{x-bf-}/\texttt{x-bifrost-}.  Body overrides are
\texttt{body.model}, \texttt{settings.body.model}, and
\texttt{settings.extraBody.model}.

Wrong-type relevant containers are refused rather than discarded.  The frame
preserves agents (including maintenance pins), sessions, separate files, other
providers, unrelated root data, exact JSON numbers, and unselected model
metadata outside those enumerated overlays.  No-op preserves bytes, inode,
mtime, and artifact count.

The checker uses a finite native-document projection rather than unrestricted
JSON.  Each owned leaf is absent, already canonical, or conflicting.  Native
registration presence is independent of the smaller certified registration
projection in $RawConfig$, so unsupported and unselected declarations can be
retained without being certified as selectable.  Presence of a selected native
registration is also independent of each selected owned leaf: a missing
$modelID$, enabled, or tools leaf is a repairable addition.  Opaque protected
sentinels cover agents (including maintenance pins), unrelated JSON, and the
exact lexical bytes of arbitrary numbers.  Distinct leaves have distinct
RFC~6901 pointers.

\begin{zed}
	OwnedLeaf ::= rootDefaultLeaf | packageLeaf | envLeaf | baseURLLeaf | \\
	\t1 apiKeyLeaf | providerCanonicalLeaf | modelRouteOverlayLeaf | \\
	\t1 variantRouteOverlayLeaf | selectedModelIDLeaf | selectedDisabledLeaf | \\
	\t1 selectedToolsLeaf \\
	LeafState ::= leafMissing | leafCanonical | leafConflict \\
	ProtectedSentinel ::= agentsSentinel | maintenancePinSentinel | \\
	\t1 sessionSentinel | separateAgentFileSentinel | unrelatedRootSentinel | \\
	\t1 otherProviderSentinel | unselectedMetadataSentinel | exactNumberSentinel
\end{zed}

\begin{axdef}
	leafPointer : OwnedLeaf \fun JSONPointer
	\where
	\forall l_1, l_2 : OwnedLeaf @
	leafPointer~l_1 = leafPointer~l_2 \implies l_1 = l_2
\end{axdef}

\begin{schema}{NativeDocument}
	leafState : OwnedLeaf \fun LeafState \\
	nativeRegistrations : \power ModelId \\
	protected : ProtectedSentinel \fun ByteString
\end{schema}

Normalization is an executable finite relation.  It assigns every raw and
native component, preserves prior fixture registrations during A-to-B switches,
establishes one coherent selected route, and frames every protected sentinel.
The changed and conflict sets are exact extensional images, not unconstrained
summaries.  Root-default replacement is changed but deliberately not a conflict;
all other conflicting owned leaves trigger the exact backup biconditional.

\begin{schema}{OpenCodeNormalization}
	\Delta RawConfig \\
	\Delta NativeDocument \\
	requested : ModelId \\
	certifiedRegistrations : \power ModelId \\
	changedPaths, conflictPaths : \power JSONPointer \\
	backupRequired : Bool
	\where
	\theta RawConfig \in Config \\
	bifrostFixture \in configured \\
	requested \in \{ sol, astra \} \\
	nativeRegistrations \subseteq \{ sol, astra, unsupported \} \\
	( opencode \in \dom registered \implies
	certifiedRegistrations = registered~opencode \land
	registered~opencode \subseteq nativeRegistrations ) \\
	( opencode \notin \dom registered \implies
	certifiedRegistrations = \emptyset ) \\
	configured' = configured \\
	routeVia' = routeVia \oplus \{ opencode \mapsto bifrostFixture \} \\
	model' = model \oplus \{ opencode \mapsto requested \} \\
	hauth' = hauth \oplus \{ opencode \mapsto vkauth \} \\
	override' = override \\
	aliases' = aliases \\
	registered' = registered \oplus
	\{ opencode \mapsto
	( ( certifiedRegistrations \cap \{ sol, astra \} ) \cup \{ requested \} ) \} \\
	shown' = shown \oplus \{ opencode \mapsto \emptyset \} \\
	nativeRegistrations' = nativeRegistrations \cup \{ requested \} \\
	protected' = protected \\
	leafState' = \{ l : OwnedLeaf @ l \mapsto leafCanonical \} \\
	changedPaths = \{ l : OwnedLeaf | leafState~l \neq leafCanonical @
	leafPointer~l \} \\
	conflictPaths = \{ l : OwnedLeaf | l \neq rootDefaultLeaf \land
	leafState~l = leafConflict @ leafPointer~l \} \\
	( backupRequired = yes \iff conflictPaths \neq \emptyset )
\end{schema}

\begin{theorem}{NormalizationEstablishesConfig}
	\forall OpenCodeNormalization @ \theta RawConfig' \in Config
\end{theorem}

\begin{theorem}{NormalizationFramesProtectedState}
	\forall OpenCodeNormalization @ protected' = protected
\end{theorem}

\begin{theorem}{NormalizationRetainsNativeRegistrations}
	\forall OpenCodeNormalization @
	nativeRegistrations \subseteq nativeRegistrations'
\end{theorem}

\begin{theorem}{NormalizationRepairsSelectedLeaves}
	\forall OpenCodeNormalization @ \\
	\t1 leafState'~selectedModelIDLeaf = leafCanonical \land \\
	\t1 leafState'~selectedDisabledLeaf = leafCanonical \land
	leafState'~selectedToolsLeaf = leafCanonical
\end{theorem}

$conflictPaths \neq \emptyset$ exactly when an existing owned routing or
registration value, other than root-default assignment alone, is replaced or
removed.  It is the backup trigger.  Additions, canonical values, an absent
selected native declaration, missing selected-registration leaves, and
root-default-only A-to-B changes require no backup.  Unselected native
declarations are retained but are not added to the certified projection.

\subsection{Closed local and live result algebra}

The public JSON schema is closed: every displayed field is required, arrays are
never null, and null is allowed only by the selected union variant.  Schema
version is integer one.  Findings sort by phase, code, and nullable field;
warnings sort by phase and code.  Human rendering JSON-quotes all externally
derived strings.  Messages never contain credentials, old/new values, full
configuration, response bodies, or raw URL-bearing errors.

\begin{zed}
	Phase ::= inputPhase | localPhase | livePhase | writePhase | reportPhase
\end{zed}

% POINTS-CHECK BEGIN types.transport type
\begin{zed}
	Transport ::= httpsTransport | loopbackHTTP | insecureHTTP
\end{zed}
% POINTS-CHECK END types.transport

\begin{zed}
	Durability ::= notApplicable | confirmed | uncertain \\
	LocalCause ::= localInvalid | localCancelled | localAbsent | localUnreadable | \\
	\t1 localMalformed | localUnmanaged | localIncoherent | localCoherent \\
	LocalVariant ::= localNotRun | localReadFailure | localNotApplicable | \\
	\t1 localIncoherentResult | localPassed \\
	LiveCause ::= liveInvalidInput | liveObservationFailed | liveRouteUnmanaged | \\
	\t1 liveLocalFailed | liveTransportNotAuthorized | liveCredentialMissing | \\
	\t1 liveCredentialInvalid | liveCancelledBeforeObservation | \\
	\t1 liveCancelledAfterLocal | liveTimeout | \\
	\t1 liveCancelled | liveUnreachable | liveRedirect | liveAuthentication | \\
	\t1 liveOtherHTTP | liveInvalidCatalog | liveModelAbsent | livePresent \\
	LiveReason ::= reasonInvalidInput | reasonObservationFailed | \\
	\t1 reasonRouteUnmanaged | reasonLocalFailed | reasonTransportNotAuthorized | \\
	\t1 reasonCredentialMissing | reasonCredentialInvalid | \\
	\t1 reasonCancelledBeforeRequest | reasonTimeout | reasonCancelled | \\
	\t1 reasonUnreachable | reasonRedirect | reasonAuthentication | reasonOtherHTTP | \\
	\t1 reasonInvalidCatalog | reasonModelAbsent \\
	LiveVariant ::= liveNotRun | liveFailed | livePassed \\
	CommandOutcome ::= unchangedOutcome | committedOutcome | refusedOutcome | \\
	\t1 failedOutcome | observedOutcome | passedOutcome \\
	FindingCode ::= eUsage | eUnsupportedHarness | eUnsupportedSource | \\
	\t1 eModelNotInFixture | eBaseURLInvalid | eInsecureHTTP | \\
	\t1 eConfigNotFound | eConfigUnreadable | eConfigFormat | eConfigShape | \\
	\t1 eRouteNotConfigured | eTransportMismatch | eCredentialReferenceMismatch | \\
	\t1 eModelNotRegistered | eRouteOverrideConflict | eParentMissing | ePathUnsafe | \\
	\t1 eWriteBusy | eBackupExists | eBackupWrite | eConfigChanged | eConfigWrite | \\
	\t1 eCleanup | eDurability | eCredentialMissing | eCredentialInvalid | \\
	\t1 eLiveRedirect | eLiveTimeout | eLiveAuth | eLiveUnreachable | eLiveHTTP | \\
	\t1 eLiveCatalog | eLiveModelAbsent | eCancelled | eReport \\
	WarningCode ::= wRouteOverwrite | wInsecureHTTP \\
	ResidualState ::= retainedResidual | deletionUnconfirmed | replacedResidual | \\
	\t1 uninspectedResidual \\
	CrashPresence ::= ownedPresent | ownedAbsent | presenceUncertain \\
	StatusVariant ::= statusObserved | statusRefused | statusFailed \\
	VerifyVariant ::= verifyPassed | verifyRefused | verifyFailed
\end{zed}

Stable machine semantics use codes, never message matching.  Phase legality and
the compact local/live cause-to-code abstraction are total finite functions.

\begin{axdef}
	allowedPhases : FindingCode \fun \power Phase \\
	liveCodes : LiveCause \fun \power FindingCode
	\where
	( \forall c : \{ eUsage, eUnsupportedHarness, eUnsupportedSource \} @
	allowedPhases~c = \{ inputPhase \} ) \\
	( \forall c : \{ eModelNotInFixture, eBaseURLInvalid \} @
	allowedPhases~c = \{ inputPhase, localPhase \} ) \\
	allowedPhases~eInsecureHTTP = \{ inputPhase, livePhase \} \\
	( \forall c : \{ eConfigNotFound, eConfigUnreadable, eConfigFormat,
	eConfigShape, eRouteNotConfigured, eTransportMismatch,
	eCredentialReferenceMismatch, eModelNotRegistered,
	eRouteOverrideConflict, eParentMissing, ePathUnsafe, eWriteBusy \} @
	allowedPhases~c = \{ localPhase \} ) \\
	allowedPhases~eBackupExists = \{ localPhase, writePhase \} \\
	( \forall c : \{ eBackupWrite, eConfigChanged, eConfigWrite, eCleanup,
	eDurability \} @ allowedPhases~c = \{ writePhase \} ) \\
	( \forall c : \{ eCredentialMissing, eCredentialInvalid, eLiveRedirect,
	eLiveTimeout, eLiveAuth, eLiveUnreachable, eLiveHTTP, eLiveCatalog,
	eLiveModelAbsent \} @ allowedPhases~c = \{ livePhase \} ) \\
	allowedPhases~eCancelled = \{ localPhase, livePhase, writePhase \} \\
	allowedPhases~eReport = \{ reportPhase \} \\
	( \forall c : \{ liveInvalidInput, liveObservationFailed,
	liveLocalFailed, livePresent \} @
	liveCodes~c = \emptyset ) \\
	liveCodes~liveRouteUnmanaged = \{ eRouteNotConfigured \} \\
	liveCodes~liveTransportNotAuthorized = \{ eInsecureHTTP \} \\
	liveCodes~liveCredentialMissing = \{ eCredentialMissing \} \\
	liveCodes~liveCredentialInvalid = \{ eCredentialInvalid \} \\
	liveCodes~liveCancelledBeforeObservation = \{ eCancelled \} \\
	liveCodes~liveCancelledAfterLocal = \{ eCancelled \} \\
	liveCodes~liveTimeout = \{ eLiveTimeout \} \\
	liveCodes~liveCancelled = \{ eCancelled \} \\
	liveCodes~liveUnreachable = \{ eLiveUnreachable \} \\
	liveCodes~liveRedirect = \{ eLiveRedirect \} \\
	liveCodes~liveAuthentication = \{ eLiveAuth \} \\
	liveCodes~liveOtherHTTP = \{ eLiveHTTP \} \\
	liveCodes~liveInvalidCatalog = \{ eLiveCatalog \} \\
	liveCodes~liveModelAbsent = \{ eLiveModelAbsent \}
\end{axdef}

\begin{axdef}
	localProducer : LocalCause \fun LocalVariant \\
	liveProducer : LiveCause \fun LiveVariant \\
	publicLiveReason : LiveCause \pfun LiveReason
	\where
	( \forall c : \{ localInvalid, localCancelled \} @
	localProducer~c = localNotRun ) \\
	( \forall c : \{ localAbsent, localUnreadable, localMalformed \} @
	localProducer~c = localReadFailure ) \\
	localProducer~localUnmanaged = localNotApplicable \\
	localProducer~localIncoherent = localIncoherentResult \\
	localProducer~localCoherent = localPassed \\
	( \forall c : \{ liveInvalidInput, liveObservationFailed,
	liveRouteUnmanaged, liveLocalFailed, liveTransportNotAuthorized,
	liveCredentialMissing, liveCredentialInvalid,
	liveCancelledBeforeObservation, liveCancelledAfterLocal \} @
	liveProducer~c = liveNotRun ) \\
	( \forall c : \{ liveTimeout, liveCancelled, liveUnreachable,
	liveRedirect, liveAuthentication, liveOtherHTTP, liveInvalidCatalog,
	liveModelAbsent \} @ liveProducer~c = liveFailed ) \\
	liveProducer~livePresent = livePassed \\
	\dom publicLiveReason = LiveCause \setminus \{ livePresent \} \\
	publicLiveReason~liveInvalidInput = reasonInvalidInput \\
	publicLiveReason~liveObservationFailed = reasonObservationFailed \\
	publicLiveReason~liveRouteUnmanaged = reasonRouteUnmanaged \\
	publicLiveReason~liveLocalFailed = reasonLocalFailed \\
	publicLiveReason~liveTransportNotAuthorized = reasonTransportNotAuthorized \\
	publicLiveReason~liveCredentialMissing = reasonCredentialMissing \\
	publicLiveReason~liveCredentialInvalid = reasonCredentialInvalid \\
	publicLiveReason~liveCancelledBeforeObservation = reasonCancelledBeforeRequest \\
	publicLiveReason~liveCancelledAfterLocal = reasonCancelledBeforeRequest \\
	publicLiveReason~liveTimeout = reasonTimeout \\
	publicLiveReason~liveCancelled = reasonCancelled \\
	publicLiveReason~liveUnreachable = reasonUnreachable \\
	publicLiveReason~liveRedirect = reasonRedirect \\
	publicLiveReason~liveAuthentication = reasonAuthentication \\
	publicLiveReason~liveOtherHTTP = reasonOtherHTTP \\
	publicLiveReason~liveInvalidCatalog = reasonInvalidCatalog \\
	publicLiveReason~liveModelAbsent = reasonModelAbsent
\end{axdef}

\begin{schema}{LocalTerminal}
	localCause : LocalCause \\
	localSelectionPresent : Bool \\
	localFindingCodes : \power FindingCode
	\where
	( localCause = localCoherent \iff localSelectionPresent = yes ) \lor
	localCause = localIncoherent \\
	( localCause = localInvalid \implies localFindingCodes = \{ eUsage \} ) \\
	( localCause = localCancelled \implies localFindingCodes = \{ eCancelled \} ) \\
	( localCause = localAbsent \implies localFindingCodes = \{ eConfigNotFound \} ) \\
	( localCause = localUnreadable \implies
	localFindingCodes = \{ eConfigUnreadable \} ) \\
	( localCause = localMalformed \implies localFindingCodes = \{ eConfigFormat \} ) \\
	( localCause \in \{ localUnmanaged, localCoherent \} \implies
	localFindingCodes = \emptyset ) \\
	( localCause = localIncoherent \implies localFindingCodes \neq \emptyset \land
	localFindingCodes \subseteq \{ eModelNotInFixture, eBaseURLInvalid,
	eConfigShape, eTransportMismatch, eCredentialReferenceMismatch,
	eModelNotRegistered, eRouteOverrideConflict \} )
\end{schema}

\begin{schema}{LocalResult}
	LocalTerminal \\
	localVariant : LocalVariant
	\where
	localVariant = localProducer~localCause
\end{schema}

\begin{theorem}{LocalProducerTotalExclusive}
	\forall LocalTerminal @
	\# \{ v : LocalVariant | v = localProducer~localCause \} = 1
\end{theorem}

The live cause is an independent terminal fact.  Its result, request flag, URL
nullability, status range, model nullability, and finding set are all derived by
one closed truth table.  Redirect is exactly 3xx; authentication is 401/403;
other HTTP excludes 2xx, 3xx, and authentication; catalog/model decisions require
2xx.  Timeout and in-flight cancellation may retain one received status.  An
early cancellation before observation has no catalog URL; cancellation after a
local pass and URL derivation retains that URL while still reporting the single
public cancelled-before-request reason.

\begin{schema}{LiveTerminal}
	liveCause : LiveCause \\
	httpStatus : \power \num
	\where
	\# httpStatus \leq 1 \\
	liveCause \in \{ liveInvalidInput, liveObservationFailed,
	liveRouteUnmanaged, liveLocalFailed, liveTransportNotAuthorized,
	liveCredentialMissing, liveCredentialInvalid,
	liveCancelledBeforeObservation, liveCancelledAfterLocal,
	liveUnreachable \} \implies
	httpStatus = \emptyset \\
	liveCause \in \{ liveTimeout, liveCancelled \} \implies
	httpStatus \subseteq 100 \upto 599 \\
	liveCause = liveRedirect \implies
	httpStatus \subseteq 300 \upto 399 \land httpStatus \neq \emptyset \\
	liveCause = liveAuthentication \implies
	httpStatus \in \{ \{ 401 \}, \{ 403 \} \} \\
	liveCause = liveOtherHTTP \implies httpStatus \neq \emptyset \land
	httpStatus \subseteq ( (100 \upto 199) \cup (400 \upto 599) )
	\setminus \{ 401, 403 \} \\
	liveCause \in \{ liveInvalidCatalog, liveModelAbsent, livePresent \}
	\implies httpStatus \neq \emptyset \land
	httpStatus \subseteq 200 \upto 299
\end{schema}

\begin{schema}{LiveResult}
	LiveTerminal \\
	liveVariant : LiveVariant \\
	requestAttempted, catalogURLPresent : Bool \\
	modelPresent : \power Bool \\
	liveResultReason : \power LiveReason \\
	liveFindingCodes : \power FindingCode
	\where
	\# modelPresent \leq 1 \\
	\# liveResultReason \leq 1 \\
	liveVariant = liveProducer~liveCause \\
	( liveCause = livePresent \iff liveResultReason = \emptyset ) \\
	( liveCause \neq livePresent \implies
	liveResultReason = \{ publicLiveReason~liveCause \} ) \\
	( requestAttempted = no \iff liveVariant = liveNotRun ) \\
	( catalogURLPresent = no \iff liveCause \in
	\{ liveInvalidInput, liveObservationFailed, liveRouteUnmanaged,
	liveLocalFailed, liveCancelledBeforeObservation \} ) \\
	( modelPresent = \{ yes \} \iff liveCause = livePresent ) \\
	( modelPresent = \{ no \} \iff liveCause = liveModelAbsent ) \\
	liveCause \notin \{ livePresent, liveModelAbsent \} \implies
	modelPresent = \emptyset \\
	liveFindingCodes = liveCodes~liveCause
\end{schema}

\begin{theorem}{LiveProducerTotalExclusive}
	\forall LiveTerminal @
	\# \{ v : LiveVariant | v = liveProducer~liveCause \} = 1
\end{theorem}

\begin{schema}{Finding}
	findingCode : FindingCode \\
	findingPhase : Phase \\
	findingField : \power JSONPointer
	\where
	\# findingField \leq 1 \\
	findingPhase \in allowedPhases~findingCode
\end{schema}

\begin{schema}{Warning}
	warningCode : WarningCode \\
	warningPhase : Phase \\
	warningConflicts : \power JSONPointer
	\where
	warningCode = wRouteOverwrite \implies warningPhase = writePhase \land \\
	\t1 warningConflicts \neq \emptyset \\
	warningCode = wInsecureHTTP \implies \\
	\t1 warningPhase \in \{ inputPhase, livePhase \} \land
	warningConflicts = \emptyset
\end{schema}

\begin{schema}{Residual}
	residualPath : Path \\
	residualState : ResidualState \\
	crashPresence : CrashPresence
	\where
	residualState = deletionUnconfirmed \implies crashPresence = presenceUncertain
\end{schema}

\begin{schema}{StatusTerminal}
	LocalResult \\
	statusReportFailed : Bool
\end{schema}

\begin{schema}{StatusObservedProducer}
	StatusTerminal
	\where
	localCause \in \{ localUnmanaged, localCoherent \} \land
	statusReportFailed = no
\end{schema}

\begin{schema}{StatusRefusedProducer}
	StatusTerminal
	\where
	localCause = localInvalid \land statusReportFailed = no
\end{schema}

\begin{schema}{StatusFailedProducer}
	StatusTerminal
	\where
	\lnot StatusObservedProducer \land \lnot StatusRefusedProducer
\end{schema}

\begin{schema}{StatusResult}
	StatusTerminal \\
	statusVariant : StatusVariant \\
	statusFindingCodes : \power FindingCode
	\where
	( statusVariant = statusObserved \iff StatusObservedProducer ) \\
	( statusVariant = statusRefused \iff StatusRefusedProducer ) \\
	( statusVariant = statusFailed \iff StatusFailedProducer ) \\
	statusReportFailed = no \implies
	statusFindingCodes = localFindingCodes \\
	statusReportFailed = yes \implies
	statusFindingCodes = localFindingCodes \cup \{ eReport \}
\end{schema}

\begin{theorem}{StatusProducerTotalExclusive}
	\forall StatusTerminal @ \\
	\# \{ v : StatusVariant |
	( v = statusObserved \land StatusObservedProducer ) \lor
	( v = statusRefused \land StatusRefusedProducer ) \lor
	( v = statusFailed \land StatusFailedProducer ) \} = 1
\end{theorem}

\begin{schema}{VerifyTerminal}
	LocalResult \\
	LiveResult \\
	verifyReportFailed : Bool
	\where
	localCause = localInvalid \iff liveCause = liveInvalidInput \\
	localCause \in \{ localAbsent, localUnreadable, localMalformed \}
	\iff liveCause = liveObservationFailed \\
	localCause = localCancelled \iff liveCause = liveCancelledBeforeObservation \\
	localCause = localUnmanaged \iff liveCause = liveRouteUnmanaged \\
	localCause = localIncoherent \iff liveCause = liveLocalFailed \\
	localCause = localCoherent \iff liveCause \notin
	\{ liveInvalidInput, liveObservationFailed, liveRouteUnmanaged,
	liveLocalFailed, liveCancelledBeforeObservation \}
\end{schema}

\begin{schema}{VerifyPassedProducer}
	VerifyTerminal
	\where
	localCause = localCoherent \land liveCause = livePresent \land
	verifyReportFailed = no
\end{schema}

\begin{schema}{VerifyRefusedProducer}
	VerifyTerminal
	\where
	localCause = localInvalid \land liveCause = liveInvalidInput \land
	verifyReportFailed = no
\end{schema}

\begin{schema}{VerifyFailedProducer}
	VerifyTerminal
	\where
	\lnot VerifyPassedProducer \land \lnot VerifyRefusedProducer
\end{schema}

\begin{schema}{VerifyResultPublic}
	VerifyTerminal \\
	verifyVariant : VerifyVariant \\
	verifyFindingCodes : \power FindingCode
	\where
	( verifyVariant = verifyPassed \iff VerifyPassedProducer ) \\
	( verifyVariant = verifyRefused \iff VerifyRefusedProducer ) \\
	( verifyVariant = verifyFailed \iff VerifyFailedProducer ) \\
	verifyReportFailed = no \implies verifyFindingCodes =
	localFindingCodes \cup liveFindingCodes \\
	verifyReportFailed = yes \implies verifyFindingCodes =
	localFindingCodes \cup liveFindingCodes \cup \{ eReport \}
\end{schema}

\begin{theorem}{VerifyProducerTotalExclusive}
	\forall VerifyTerminal @ \\
	\# \{ v : VerifyVariant |
	( v = verifyPassed \land VerifyPassedProducer ) \lor
	( v = verifyRefused \land VerifyRefusedProducer ) \lor
	( v = verifyFailed \land VerifyFailedProducer ) \} = 1
\end{theorem}

The local classification is total: invalid input/cancellation before observation
is not-run; missing/unreadable/malformed input fails; a supported document whose
default does not select Bifrost is unmanaged; a selected but invalid Bifrost
route is incoherent; only the exact managed projection is coherent.  Status
succeeds for unmanaged or coherent state and never performs live work.

Verify's ordered gates are input, observation, managed-route classification,
local coherence, HTTP authorization, credential validation, and request.  The
first blocked gate determines the live not-run reason.  Local failure never
reads the credential or attempts a request.  A passed verify requires both
local and live pass; not-run can never pass.

Live verification resolves \texttt{BIFROST\_VK} in memory, validates it as a
header value, and makes exactly one \texttt{GET R/v1/models} with bearer auth.
The client follows no redirects, retries nothing, has a ten-second total
deadline and four-MiB decoded-body limit, ignores environment proxies, and uses
a fresh transport without implicit retry.  A valid response is a successful
object containing a \texttt{data} array of objects with string \texttt{id}; the
selected token must occur exactly.  Authentication, redirect, timeout,
unreachable, other HTTP, invalid catalog, and model absence stay distinct.

\section{Completed route results and exhaustive producers}

An internal terminal event has an immutable class, code, and phase.  Class is
producer metadata, not a public JSON field.  Refusals are policy/format/shape
rejections before writing; execution, cancellation, cleanup, durability, and
reporting errors are operational even when no mutation happened.

\begin{zed}
	EventClass ::= refusalClass | operationalClass
\end{zed}

% POINTS-CHECK BEGIN definitions.events definition
\begin{zed}
	RouteEvent ::= iUsage | iHarness | iSource | iModel | iURL | iHTTP | \\
	\t1 lParent | lUnsafe | lBusy | lUnreadable | lFormat | lShape | lBackupExists | \\
	\t1 wCandidate | wConfigStage | wBackupStage | wBackupCollision | wBackupPublish | \\
	\t1 wBackupCleanup | wChanged | wAbsentOccupied | wConfigPublish | dConfigSync | \\
	\t1 cUnlink | cReplacement | cInspect | cDeleteSync | cRelease | \\
	\t1 xBeforeInput | xBeforeWrite | xDuringWrite | xAfterCommit | \\
	\t1 oReportBefore | oReportAfter
\end{zed}
% POINTS-CHECK END definitions.events

\begin{zed}
	RouteVariant ::= routeUnchanged | routeCommitted | routeRefused | \\
	\t1 routeFailedBefore | routeFailedAfter
\end{zed}

\begin{axdef}
	refusalEvents, operationalEvents : \power RouteEvent \\
	eventClass : RouteEvent \fun EventClass \\
	eventPhase : RouteEvent \fun Phase \\
	eventCode : RouteEvent \fun FindingCode
	\where
	refusalEvents = \{ iUsage, iHarness, iSource, iModel, iURL, iHTTP,
	lParent, lUnsafe, lBusy, lFormat, lShape, lBackupExists \} \\
	operationalEvents = RouteEvent \setminus refusalEvents \\
	( \forall e : refusalEvents @ eventClass~e = refusalClass ) \\
	( \forall e : operationalEvents @ eventClass~e = operationalClass ) \\
	eventPhase~iUsage = inputPhase \\
	eventPhase~iHarness = inputPhase \\
	eventPhase~iSource = inputPhase \\
	eventPhase~iModel = inputPhase \\
	eventPhase~iURL = inputPhase \\
	eventPhase~iHTTP = inputPhase \\
	( \forall e : \{ lParent, lUnsafe, lBusy, lUnreadable, lFormat, lShape,
	lBackupExists, xBeforeInput, xBeforeWrite \} @ eventPhase~e = localPhase ) \\
	( \forall e : \{ wCandidate, wConfigStage, wBackupStage, wBackupCollision,
	wBackupPublish, wBackupCleanup, wChanged, wAbsentOccupied, wConfigPublish,
	dConfigSync, cUnlink, cReplacement, cInspect, cDeleteSync, cRelease,
	xDuringWrite, xAfterCommit \} @ eventPhase~e = writePhase ) \\
	eventPhase~oReportBefore = reportPhase \\
	eventPhase~oReportAfter = reportPhase \\
	eventCode~iUsage = eUsage \\
	eventCode~iHarness = eUnsupportedHarness \\
	eventCode~iSource = eUnsupportedSource \\
	eventCode~iModel = eModelNotInFixture \\
	eventCode~iURL = eBaseURLInvalid \\
	eventCode~iHTTP = eInsecureHTTP \\
	eventCode~lParent = eParentMissing \\
	eventCode~lUnsafe = ePathUnsafe \\
	eventCode~lBusy = eWriteBusy \\
	eventCode~lUnreadable = eConfigUnreadable \\
	eventCode~lFormat = eConfigFormat \\
	eventCode~lShape = eConfigShape \\
	eventCode~lBackupExists = eBackupExists \\
	( \forall e : \{ wCandidate, wConfigStage, wConfigPublish \} @
	eventCode~e = eConfigWrite ) \\
	( \forall e : \{ wBackupStage, wBackupPublish \} @
	eventCode~e = eBackupWrite ) \\
	eventCode~wBackupCollision = eBackupExists \\
	eventCode~wChanged = eConfigChanged \\
	eventCode~wAbsentOccupied = eConfigChanged \\
	( \forall e : \{ wBackupCleanup, cUnlink, cReplacement, cInspect,
	cDeleteSync, cRelease \} @ eventCode~e = eCleanup ) \\
	eventCode~dConfigSync = eDurability \\
	( \forall e : \{ xBeforeInput, xBeforeWrite, xDuringWrite, xAfterCommit \} @
	eventCode~e = eCancelled ) \\
	eventCode~oReportBefore = eReport \\
	eventCode~oReportAfter = eReport
\end{axdef}

The event-to-code/phase table is normative:

\begin{center}
	\begin{tabular}{lll}
		I1--I6 & usage/harness/source/model/URL/insecure HTTP      & input, refusal, null selection \\
		L1--L3 & parent/path/lock contention                       & local, refusal, selected       \\
		L4     & metadata, lock, snapshot, read or close I/O       & local unreadable, operational  \\
		L5--L7 & strict format, shape, early backup collision      & local, refusal                 \\
		W1--W3 & candidate/config-stage/backup-stage failure       & write, operational             \\
		W4--W6 & late collision/publication/backup cleanup         & write, operational             \\
		W7--W9 & revalidation/no-replace/config publication        & write, operational             \\
		D1     & config publication sync                           & write durability, operational  \\
		C1--C5 & unlink/replacement/inspection/delete-sync/release & write cleanup, operational     \\
		X1--X4 & cancellation at named boundaries                  & local or write, operational    \\
		O1--O2 & result preparation or delivery                    & report, operational
	\end{tabular}
\end{center}

The corresponding public stable codes are, in order,
\texttt{E\_USAGE}, \texttt{E\_UNSUPPORTED\_HARNESS},
\texttt{E\_UNSUPPORTED\_SOURCE}, \texttt{E\_MODEL\_NOT\_IN\_FIXTURE},
\texttt{E\_BASE\_URL\_INVALID}, \texttt{E\_INSECURE\_HTTP},
\texttt{E\_PARENT\_MISSING}, \texttt{E\_PATH\_UNSAFE},
\texttt{E\_WRITE\_BUSY}, \texttt{E\_CONFIG\_UNREADABLE},
\texttt{E\_CONFIG\_FORMAT}, \texttt{E\_CONFIG\_SHAPE},
\texttt{E\_BACKUP\_EXISTS}, \texttt{E\_CONFIG\_WRITE},
\texttt{E\_BACKUP\_WRITE}, \texttt{E\_CONFIG\_CHANGED},
\texttt{E\_DURABILITY}, \texttt{E\_CLEANUP},
\texttt{E\_CANCELLED}, and \texttt{E\_REPORT}.  Read-only codes have no route
producer.  In particular old incoherence is repaired, not emitted by route.

Selection becomes non-null exactly when the complete input gate succeeds and is
never erased.  The gate validates syntax, supported combination, exact fixture,
absolute lexical config path, root URL, and HTTP authorization.

\begin{schema}{RouteTerminal}
	events : \power RouteEvent \\
	inputAccepted, committed, noOp, configSyncAck : Bool \\
	namespaceMutated : Bool
	\where
	( events = \emptyset \implies \\
	\t1 ( committed = yes \land noOp = no ) \lor
		( committed = no \land noOp = yes ) ) \\
	committed = yes \implies inputAccepted = yes \\
	noOp = yes \implies inputAccepted = yes \land committed = no \\
	events \cap \{ iUsage, iHarness, iSource, iModel, iURL, iHTTP,
	xBeforeInput \} \neq \emptyset \implies inputAccepted = no \\
	events \cap \{ lParent, lUnsafe, lBusy, lUnreadable, lFormat, lShape,
	lBackupExists, wCandidate, wConfigStage, wBackupStage,
	wBackupCollision, wBackupPublish, wBackupCleanup, wChanged,
	wAbsentOccupied, wConfigPublish, dConfigSync, cUnlink, cReplacement,
	cInspect, cDeleteSync, cRelease, xBeforeWrite, xDuringWrite,
	xAfterCommit \} \neq \emptyset \implies inputAccepted = yes \\
	events \cap \{ xAfterCommit, oReportAfter, dConfigSync \} \neq \emptyset
	\implies committed = yes \\
	oReportBefore \in events \implies committed = no \\
	noOp = yes \land events \neq \emptyset \implies
	events \cap operationalEvents \neq \emptyset \\
	configSyncAck = yes \implies committed = yes \\
	committed = yes \land configSyncAck = no \implies dConfigSync \in events \\
	dConfigSync \in events \implies committed = yes \\
	( events \neq \emptyset \land events \subseteq refusalEvents \implies \\
	\t1 committed = no \land namespaceMutated = no )
\end{schema}

\begin{schema}{RouteUnchangedProducer}
	RouteTerminal
	\where
	committed = no \land events = \emptyset \land noOp = yes
\end{schema}

\begin{schema}{RouteCommittedProducer}
	RouteTerminal
	\where
	committed = yes \land events = \emptyset
\end{schema}

\begin{schema}{RouteRefusedProducer}
	RouteTerminal
	\where
	committed = no \land events \neq \emptyset \land
	events \cap operationalEvents = \emptyset
\end{schema}

\begin{schema}{RouteFailedBeforeProducer}
	RouteTerminal
	\where
	committed = no \land events \neq \emptyset \land
	events \cap operationalEvents \neq \emptyset
\end{schema}

\begin{schema}{RouteFailedAfterProducer}
	RouteTerminal
	\where
	committed = yes \land events \neq \emptyset
\end{schema}

\begin{schema}{CompletedRoute}
	RouteTerminal \\
	variant : RouteVariant \\
	durability : Durability
	\where
	( variant = routeUnchanged \iff RouteUnchangedProducer ) \\
	( variant = routeCommitted \iff RouteCommittedProducer ) \\
	( variant = routeRefused \iff RouteRefusedProducer ) \\
	( variant = routeFailedBefore \iff RouteFailedBeforeProducer ) \\
	( variant = routeFailedAfter \iff RouteFailedAfterProducer ) \\
	( committed = no \iff durability = notApplicable ) \\
	( committed = yes \land configSyncAck = yes \iff durability = confirmed ) \\
	( committed = yes \land configSyncAck = no \iff durability = uncertain )
\end{schema}

\begin{schema}{RoutePublicResult}
	CompletedRoute \\
	selectionPresent, conflictsPresent : Bool \\
	backupPaths : \power Path \\
	warningCodes : \power WarningCode \\
	findingCodes : \power FindingCode \\
	residualPaths : \power Path
	\where
	\# backupPaths \leq 1 \\
	( selectionPresent = yes \iff inputAccepted = yes ) \\
	findingCodes = eventCode \limg events \rimg \\
	( conflictsPresent = yes \land committed = yes \iff
	wRouteOverwrite \in warningCodes ) \\
	variant = routeUnchanged \implies backupPaths = \emptyset \land \\
	\t1 findingCodes = \emptyset \land residualPaths = \emptyset \\
	variant = routeCommitted \implies findingCodes = \emptyset \land \\
	\t1 residualPaths = \emptyset \\
	variant = routeRefused \implies backupPaths = \emptyset \land \\
	\t1 residualPaths = \emptyset \land namespaceMutated = no \\
	committed = yes \land conflictsPresent = yes \implies backupPaths \neq \emptyset \\
	residualPaths \neq \emptyset \implies eCleanup \in findingCodes
\end{schema}

\begin{theorem}{RouteProducerTotalExclusive}
	\forall RouteTerminal @ \\
	\t1 \# \{ v : RouteVariant |
	\t2 ( v = routeUnchanged \land RouteUnchangedProducer ) \lor
	\t2 ( v = routeCommitted \land RouteCommittedProducer ) \lor
	\t2 ( v = routeRefused \land RouteRefusedProducer ) \lor
	\t2 ( v = routeFailedBefore \land RouteFailedBeforeProducer ) \lor
	\t2 ( v = routeFailedAfter \land RouteFailedAfterProducer ) \} = 1
\end{theorem}

\begin{theorem}{RouteDiscriminatorExistsUnique}
	\forall RouteTerminal @ \exists_1 v : RouteVariant @ \\
	\t1 ( v = routeUnchanged \land RouteUnchangedProducer ) \lor
	\t1 ( v = routeCommitted \land RouteCommittedProducer ) \lor
	\t1 ( v = routeRefused \land RouteRefusedProducer ) \lor
	\t1 ( v = routeFailedBefore \land RouteFailedBeforeProducer ) \lor
	\t1 ( v = routeFailedAfter \land RouteFailedAfterProducer )
\end{theorem}

\begin{theorem}{MissingConfigSyncProducesDurabilityEvent}
	\forall RouteTerminal @ committed = yes \land configSyncAck = no
	\implies dConfigSync \in events
\end{theorem}

\begin{theorem}{HistoricalConfigSyncFailureCanBeLaterConfirmed}
	\exists CompletedRoute @ dConfigSync \in events \land
	configSyncAck = yes \land durability = confirmed \land
	variant = routeFailedAfter
\end{theorem}

The biconditionals are the producer, not loose constructor shapes.  They prove
totality by partitioning first on commit, then findings, then operational event,
and prove pairwise exclusivity directly.  Every named event is in exactly one
event class.  Reporting failure adds an operational event and therefore cannot
retain a success/refusal variant.  A syntax refusal delivered successfully exits
2; successful route exits 0; all other route results, including reporting
failure after usage error, exit 1.

The route result carries schema version, command, explicit-file scope, context,
nullable selection and owned backup path, warnings, findings, and residuals.
Unchanged is uncommitted/not-applicable with no findings or residue.  Committed
is committed/confirmed with no findings or residue.  Refused has no namespace
mutation, backup, or residue.  Operational pre-commit failure may have no
mutation or may retain owned residue.  Post-commit failure reports confirmed or
uncertain durability truthfully.  Selection describes validated request, not
installation success.

Exactly one \texttt{W\_ROUTE\_OVERWRITE} is emitted iff a conflicting candidate
commits; it includes the exact backup path and sorted conflict pointers and
survives later failure.  \texttt{W\_INSECURE\_HTTP} occurs at most once after
the relevant authorization gate and survives later failure.  Status emits no
warnings.  A committed backup is a retained result, never residue.

Residuals are a closed union: retained owned name after failed unlink,
deletion-unconfirmed after visible unlink without directory-sync evidence,
replacement identity left untouched, or uninspected path.  Crash presence is
owned-present, owned-absent, or uncertain.  Every residue entails cleanup error;
success has none.  A detected anomaly remains reportable even if later sync
proves the owned identity absent.

\section{Filesystem publication and paired crash model}

The executable Increment~1 transaction library appears after the following
historical model.  The historical three-slot draft is retained temporarily as
review context but is not part of the active specification or any generated
checker projection.

\begin{verbatim}
	The filesystem model uses finite carriers under model-checker bounds.  A file
	identity has one image even when multiple paths hard-link it.

	Only strict JSON objects are accepted: no comments, trailing commas, duplicate
	decoded keys, multiple documents, non-whitespace trailing data, non-object root,
	invalid UTF-8, or malformed strings.  Original bytes are retained separately
	for backup.  The explicit immediate parent must exist, be effective-user-owned,
	and not group/other writable; ancestors are user/root controlled and not
	untrusted-writable.  Descriptor-relative access rejects symlink traversal.
	Existing targets are regular, effective-user-owned, link-count one, owner
	read/write, non-executable, without special or group/other-write bits, and with
	supported mode-equivalent access metadata.  Incompatible ACL/default-ACL effects
	fail closed.  New config mode is 0600; replacement and backup preserve the
	supported original mode.  Timestamp, inode, xattr, and arbitrary ACL
	preservation are not promised.

	Route takes a nonblocking exclusive parent-directory advisory lock from snapshot
	through cleanup; readers share-lock only while snapshotting.  Revalidation before
	commit compares target existence, identity, exact bytes, ownership, link count,
	mode, and supported access metadata.  Mismatch retains the external edit and
	produces config-changed.  The guarantee excludes non-cooperating mutation in the
	final revalidate/publication window and same-UID hostile interference: rename is
	not compare-and-swap.

	For config $P$, the final backup is exactly $P$ plus
	\texttt{.points.bak}.  Sibling stages are
	\texttt{.<base>.points-config.<nonce>.tmp} and
	\texttt{.<base>.points-backup.<nonce>.tmp}; a nonce is 128 random bits rendered
	as 32 lower-case hexadecimal digits, with at most 32 exclusive-create attempts.
	Stages may be partial, but publication requires complete bytes, established and
	verified mode/access metadata, successful file sync, and successful close.

	A required backup is staged from the exact snapshot, exclusively hard-linked at
	the final backup path, parent-synced, then has its stage name unlinked and that
	deletion parent-synced.  Only then may config commit.  Anything at the final
	backup path blocks this sequence; nothing there is followed, reused, rotated,
	truncated, or replaced.  Existing config commits by same-directory atomic
	rename.  Initially absent config commits by exclusive hard link.  The successful
	publication syscall sets committed immediately; subsequent parent sync,
	config-stage cleanup, deletion sync, and reporting are distinct outcomes.

	% historical \begin{zed}
		EntryRef ::= file \ldata Identity \rdata | symlinkRef | directoryRef | specialRef \\
		Namespace == Path \pfun EntryRef \\
		CleanupObservation ::= noCleanupNote | unlinkFailed | identityChanged | inspectFailed \\
		BackupProgress ::= backupNone | backupStaged | backupLinked | linkSynced | \\
		\t1 stageRemoved | backupReady \\
		PendingCount ::= pendingZero | pendingOne | pendingTwo | pendingThree \\
		TransactionPhase ::= validating | staging | publishingBackup | revalidating | \\
		\t1 publishingConfig | confirming | cleaning | reporting | aborting | finished \\
		Life ::= running | ended \\
		ReportState ::= reportPending | reportDelivered | reportFailed | processDead
	\end{zed}

	% historical \begin{axdef}
		deltaRemove : Delta \fun \power Path \\
		deltaPut : Delta \fun Namespace \\
		applyDelta : Namespace \cross Delta \fun Namespace \\
		refs : Namespace \fun \power Identity \\
		emptyBytes : ByteString \\
		bytePrefix : ByteString \rel ByteString \\
		imageBytes : FileImage \fun ByteString \\
		imageMode : FileImage \fun Mode \\
		imageOwner : FileImage \fun Owner \\
		imageState : FileImage \fun ContentState
		\where
		\forall n : Namespace; d : Delta @ \\
		\t1 applyDelta~(n,d) = (deltaRemove~d \ndres n) \oplus deltaPut~d \\
		\forall n : Namespace @ refs~n = \\
		\t1 \{ id : Identity | \exists p : \dom n @ n~p = file~id \} \\
		\forall b : ByteString @ emptyBytes \mapsto b \in bytePrefix \\
		\forall b : ByteString @ b \mapsto b \in bytePrefix \\
		\forall a, b : ByteString @
		a \mapsto b \in bytePrefix \land b \mapsto a \in bytePrefix \implies a = b \\
		\forall a, b, c : ByteString @
		a \mapsto b \in bytePrefix \land b \mapsto c \in bytePrefix
		\implies a \mapsto c \in bytePrefix \\
		\forall b : ByteString; m : Mode; o : Owner; c : ContentState @ \\
		\t1 \exists_1 i : FileImage @ imageBytes~i = b \land imageMode~i = m \land \\
		\t2 imageOwner~i = o \land imageState~i = c \\
		\forall i, j : FileImage @ i = j \iff \\
		\t1 imageBytes~i = imageBytes~j \land imageMode~i = imageMode~j \land \\
		\t1 imageOwner~i = imageOwner~j \land imageState~i = imageState~j
	\end{axdef}

	$FileImage$ is an initial finite carrier because ProB rejects the equivalent
	late product/record declarations in this document.  The four total selectors,
	the unique existence of one atom for every field tuple, and selector
	extensionality make it a lossless finite record abstraction: bytes, mode, owner,
	and content state can neither be omitted nor aliased.  $bytePrefix$ is the
	bounded actual-prefix order supplied by a checker valuation; $emptyBytes$ is its
	least element.

	% historical \begin{schema}{Transaction}
		visibleNS, durableNS : Namespace \\
		pendingCount : PendingCount \\
		pendingDeltaOne, pendingDeltaTwo, pendingDeltaThree : Delta \\
		visibleImages, syncedImages, intendedImages, snapshotImages :
		Identity \pfun FileImage \\
		dirtyChoices : Identity \pfun \power FileImage \\
		stageIDs, sealedIDs, closeOK : \power Identity \\
		owned, cleanupDue, cleanupPending, retainedResults : Path \pfun Identity \\
		cleanupNotes : Path \pfun CleanupObservation \\
		published : \power ( Path \cross Identity ) \\
		publicationImages : Identity \pfun FileImage \\
		configPath, backupPath, configStage, backupStage : Path \\
		candidateID : Identity \\
		backupIDs : \power Identity \\
		expectedConfig : \power Identity \\
		conflicts : \power JSONPointer \\
		phase : TransactionPhase \\
		committed, configSyncAck, cancelled, lockHeld : Bool \\
		backupProgress : BackupProgress \\
		life : Life \\
		reportState : ReportState
		\where
		\# \{ configPath, backupPath, configStage, backupStage \} = 4 \\
		\# expectedConfig \leq 1 \\
		candidateID \notin backupIDs \\
		cleanupDue \subseteq owned \\
		cleanupPending \subseteq cleanupDue \\
		retainedResults \subseteq owned \\
		\dom cleanupDue \cap \dom retainedResults = \emptyset \\
		\dom syncedImages \cap \dom dirtyChoices = \emptyset \\
		( \forall choices : \ran dirtyChoices @ choices \neq \emptyset ) \\
		( \forall id : \dom intendedImages @
		imageState~(intendedImages~id) = complete ) \\
		( \forall id : \dom dirtyChoices @ \\
		\t1 id \in \dom visibleImages \cap \dom intendedImages \land \\
		\t1 imageBytes~(visibleImages~id) \mapsto
		imageBytes~(intendedImages~id) \in bytePrefix \land \\
		\t1 imageMode~(visibleImages~id) = imageMode~(intendedImages~id) \land \\
		\t1 imageOwner~(visibleImages~id) = imageOwner~(intendedImages~id) \land \\
		\t1 ( imageState~(visibleImages~id) = complete \iff \\
		\t2 imageBytes~(visibleImages~id) = imageBytes~(intendedImages~id) \land \\
		\t2 imageMode~(visibleImages~id) = imageMode~(intendedImages~id) \land
		imageOwner~(visibleImages~id) = imageOwner~(intendedImages~id) ) \land \\
		\t1 dirtyChoices~id = \{ i : FileImage | \\
		\t2 imageBytes~i \mapsto imageBytes~(visibleImages~id) \in bytePrefix \land \\
		\t2 imageMode~i = imageMode~(intendedImages~id) \land \\
		\t2 imageOwner~i = imageOwner~(intendedImages~id) \land \\
		\t2 ( imageState~i = complete \iff \\
		\t3 imageBytes~i = imageBytes~(intendedImages~id) \land
		imageMode~i = imageMode~(intendedImages~id) \land \\
		\t3 imageOwner~i = imageOwner~(intendedImages~id) ) @ i \} ) \\
		refs~visibleNS \cup refs~durableNS
		\subseteq \dom syncedImages \cup \dom dirtyChoices \\
		pendingCount \in \{ pendingOne, pendingTwo, pendingThree \} \implies
		refs~(deltaPut~pendingDeltaOne) \subseteq
		\dom syncedImages \cup \dom dirtyChoices \\
		pendingCount \in \{ pendingTwo, pendingThree \} \implies
		refs~(deltaPut~pendingDeltaTwo) \subseteq
		\dom syncedImages \cup \dom dirtyChoices \\
		pendingCount = pendingThree \implies
		refs~(deltaPut~pendingDeltaThree) \subseteq
		\dom syncedImages \cup \dom dirtyChoices \\
		( \forall p : Path; id : Identity | (p,id) \in published @ \\
		\t1 id \in \dom syncedImages \cap \dom publicationImages \land \\
		\t1 syncedImages~id = publicationImages~id \land \\
		\t1 imageState~(publicationImages~id) = complete ) \\
		backupIDs \subseteq \dom snapshotImages \\
		( \forall id : backupIDs \cap \dom publicationImages @ \\
		\t1 publicationImages~id = snapshotImages~id ) \\
		( \forall id : \dom publicationImages @
		imageState~(publicationImages~id) = complete ) \\
		life = running \land committed = yes \land conflicts \neq \emptyset \implies \\
		\t1 backupPath \in \dom retainedResults \land backupPath \notin \dom cleanupDue
	\end{schema}

	The checker never quantifies over every sequence or its powerset.  At most three
	pending namespace deltas are represented by explicit finite slots.  The
	following relation appends one delta and has no recursive definition.

	% historical \begin{schema}{AppendPendingDelta}
		pendingCount, pendingCount' : PendingCount \\
		pendingDeltaOne, pendingDeltaTwo, pendingDeltaThree : Delta \\
		pendingDeltaOne', pendingDeltaTwo', pendingDeltaThree' : Delta \\
		delta? : Delta
		\where
		( pendingCount = pendingZero \land pendingCount' = pendingOne \land
		pendingDeltaOne' = delta? \land
		pendingDeltaTwo' = pendingDeltaTwo \land
		pendingDeltaThree' = pendingDeltaThree ) \lor \\
		( pendingCount = pendingOne \land pendingCount' = pendingTwo \land
		pendingDeltaOne' = pendingDeltaOne \land
		pendingDeltaTwo' = delta? \land
		pendingDeltaThree' = pendingDeltaThree ) \lor \\
		( pendingCount = pendingTwo \land pendingCount' = pendingThree \land
		pendingDeltaOne' = pendingDeltaOne \land
		pendingDeltaTwo' = pendingDeltaTwo \land
		pendingDeltaThree' = delta? )
	\end{schema}

	Frame schemas make omitted effects impossible.  Every transition below includes
	the frames for untouched groups and explicitly assigns every component in each
	changed group.

	% historical \begin{schema}{TxnStaticFrame}
		\Delta Transaction
		\where
		configPath' = configPath \land backupPath' = backupPath \land
		configStage' = configStage \land backupStage' = backupStage \\
		candidateID' = candidateID \land backupIDs' = backupIDs \land
		expectedConfig' = expectedConfig \land conflicts' = conflicts
	\end{schema}

	% historical \begin{schema}{TxnNamespaceFrame}
		\Delta Transaction
		\where
		visibleNS' = visibleNS \land durableNS' = durableNS \land
		pendingCount' = pendingCount \\
		pendingDeltaOne' = pendingDeltaOne \land
		pendingDeltaTwo' = pendingDeltaTwo \land
		pendingDeltaThree' = pendingDeltaThree
	\end{schema}

	% historical \begin{schema}{TxnImageFrame}
		\Delta Transaction
		\where
		visibleImages' = visibleImages \land syncedImages' = syncedImages \land
		intendedImages' = intendedImages \land snapshotImages' = snapshotImages \\
		dirtyChoices' = dirtyChoices \land stageIDs' = stageIDs \land
		sealedIDs' = sealedIDs \land closeOK' = closeOK \\
		published' = published \land publicationImages' = publicationImages
	\end{schema}

	% historical \begin{schema}{TxnOwnershipFrame}
		\Delta Transaction
		\where
		owned' = owned \land cleanupDue' = cleanupDue \land
		cleanupPending' = cleanupPending \land
		retainedResults' = retainedResults \land cleanupNotes' = cleanupNotes
	\end{schema}

	% historical \begin{schema}{TxnControlFrame}
		\Delta Transaction
		\where
		phase' = phase \land committed' = committed \land
		configSyncAck' = configSyncAck \land cancelled' = cancelled \land
		lockHeld' = lockHeld \land backupProgress' = backupProgress \land
		life' = life \land reportState' = reportState
	\end{schema}

	% historical \begin{schema}{CreateStage}
		\Delta Transaction \\
		TxnStaticFrame \\
		stage? : Path; id? : Identity; initialImage? : FileImage; delta? : Delta
		\where
		life = running \land cancelled = no \land
		phase \in \{ validating, staging \} \land lockHeld = yes \\
		stage? \in \{ configStage, backupStage \} \land
		stage? \notin \dom visibleNS \cup \dom owned \\
		( stage? = configStage \implies id? = candidateID ) \\
		( stage? = backupStage \implies id? \in backupIDs ) \\
		id? \in \dom intendedImages \\
		id? \notin refs~visibleNS \cup refs~durableNS \cup \ran owned \cup
		\dom visibleImages \cup \dom syncedImages \cup \dom dirtyChoices \cup
		stageIDs \cup \dom publicationImages \\
		imageBytes~initialImage? = emptyBytes \land
		imageMode~initialImage? = imageMode~(intendedImages~id?) \land \\
		imageOwner~initialImage? = imageOwner~(intendedImages~id?) \land
		( imageState~initialImage? = complete \iff \\
		\t1 imageBytes~initialImage? = imageBytes~(intendedImages~id?) \land
		imageMode~initialImage? = imageMode~(intendedImages~id?) \land \\
		\t1 imageOwner~initialImage? = imageOwner~(intendedImages~id?) ) \\
		( id? \in backupIDs \implies id? \in \dom snapshotImages \land
		intendedImages~id? = snapshotImages~id? ) \\
		deltaRemove~delta? = \emptyset \land
		deltaPut~delta? = \{ stage? \mapsto file~id? \} \\
		AppendPendingDelta \\
		visibleNS' = applyDelta~(visibleNS,delta?) \land durableNS' = durableNS \\
		visibleImages' = visibleImages \oplus \{ id? \mapsto initialImage? \} \\
		syncedImages' = syncedImages \land intendedImages' = intendedImages \land
		snapshotImages' = snapshotImages \\
		dirtyChoices' = dirtyChoices \oplus \{ id? \mapsto
		\{ i : FileImage | \\
		\t1 imageBytes~i \mapsto imageBytes~initialImage? \in bytePrefix \land \\
		\t1 imageMode~i = imageMode~(intendedImages~id?) \land \\
		\t1 imageOwner~i = imageOwner~(intendedImages~id?) \land \\
		\t1 ( imageState~i = complete \iff \\
		\t2 imageBytes~i = imageBytes~(intendedImages~id?) \land
		imageMode~i = imageMode~(intendedImages~id?) \land \\
		\t2 imageOwner~i = imageOwner~(intendedImages~id?) ) @ i \} \} \\
		stageIDs' = stageIDs \cup \{ id? \} \land sealedIDs' = sealedIDs \land
		closeOK' = closeOK \land published' = published \land
		publicationImages' = publicationImages \\
		owned' = owned \oplus \{ stage? \mapsto id? \} \land
		cleanupDue' = cleanupDue \oplus \{ stage? \mapsto id? \} \\
		cleanupPending' = cleanupPending \land retainedResults' = retainedResults \land
		cleanupNotes' = cleanupNotes \\
		phase' = staging \land committed' = committed \land
		configSyncAck' = configSyncAck \land cancelled' = cancelled \land
		lockHeld' = lockHeld \land backupProgress' = backupProgress \land
		life' = life \land reportState' = reportState
	\end{schema}

	% historical \begin{schema}{WriteStage}
		\Delta Transaction \\
		TxnStaticFrame \\
		TxnNamespaceFrame \\
		TxnOwnershipFrame \\
		TxnControlFrame \\
		id? : Identity; image? : FileImage
		\where
		life = running \land cancelled = no \land phase = staging \land
		lockHeld = yes \\
		id? \in ( stageIDs \setminus sealedIDs ) \cap \dom dirtyChoices \cap
		\dom intendedImages \cap \dom visibleImages \\
		imageBytes~(visibleImages~id?) \mapsto imageBytes~image? \in bytePrefix \land \\
		imageBytes~image? \mapsto imageBytes~(intendedImages~id?) \in bytePrefix \\
		imageMode~image? = imageMode~(intendedImages~id?) \land
		imageOwner~image? = imageOwner~(intendedImages~id?) \\
		( imageState~image? = complete \iff \\
		\t1 imageBytes~image? = imageBytes~(intendedImages~id?) \land
		imageMode~image? = imageMode~(intendedImages~id?) \land \\
		\t1 imageOwner~image? = imageOwner~(intendedImages~id?) ) \\
		visibleImages' = visibleImages \oplus \{ id? \mapsto image? \} \\
		syncedImages' = syncedImages \land intendedImages' = intendedImages \land
		snapshotImages' = snapshotImages \\
		dirtyChoices' = dirtyChoices \oplus \{ id? \mapsto
		\{ i : FileImage | \\
		\t1 imageBytes~i \mapsto imageBytes~image? \in bytePrefix \land \\
		\t1 imageMode~i = imageMode~(intendedImages~id?) \land \\
		\t1 imageOwner~i = imageOwner~(intendedImages~id?) \land \\
		\t1 ( imageState~i = complete \iff \\
		\t2 imageBytes~i = imageBytes~(intendedImages~id?) \land
		imageMode~i = imageMode~(intendedImages~id?) \land \\
		\t2 imageOwner~i = imageOwner~(intendedImages~id?) ) @ i \} \} \\
		stageIDs' = stageIDs \land sealedIDs' = sealedIDs \land closeOK' = closeOK \\
		published' = published \land publicationImages' = publicationImages
	\end{schema}

	% historical \begin{schema}{FileSyncStage}
		\Delta Transaction \\
		TxnStaticFrame \\
		TxnNamespaceFrame \\
		TxnOwnershipFrame \\
		TxnControlFrame \\
		id? : Identity
		\where
		life = running \land cancelled = no \land phase = staging \land
		lockHeld = yes \\
		id? \in stageIDs \cap \dom visibleImages \cap \dom intendedImages \\
		visibleImages~id? = intendedImages~id? \land
		imageState~(intendedImages~id?) = complete \\
		visibleImages' = visibleImages \land
		syncedImages' = syncedImages \oplus \{ id? \mapsto intendedImages~id? \} \\
		intendedImages' = intendedImages \land snapshotImages' = snapshotImages \\
		dirtyChoices' = \{ id? \} \ndres dirtyChoices \\
		stageIDs' = stageIDs \land sealedIDs' = sealedIDs \cup \{ id? \} \land
		closeOK' = closeOK \land published' = published \land
		publicationImages' = publicationImages
	\end{schema}

	% historical \begin{schema}{CloseStage}
		\Delta Transaction \\
		TxnStaticFrame \\
		TxnNamespaceFrame \\
		TxnOwnershipFrame \\
		id? : Identity
		\where
		life = running \land cancelled = no \land phase = staging \land
		lockHeld = yes \\
		id? \in sealedIDs \cap stageIDs \\
		visibleImages' = visibleImages \land syncedImages' = syncedImages \land
		intendedImages' = intendedImages \land snapshotImages' = snapshotImages \land
		dirtyChoices' = dirtyChoices \land stageIDs' = stageIDs \land
		sealedIDs' = sealedIDs \land closeOK' = closeOK \cup \{ id? \} \\
		published' = published \land publicationImages' = publicationImages \\
		phase' = phase \land committed' = committed \land
		configSyncAck' = configSyncAck \land cancelled' = cancelled \land
		lockHeld' = lockHeld \\
		( id? \in backupIDs \implies backupProgress' = backupStaged ) \\
		( id? \notin backupIDs \implies backupProgress' = backupProgress ) \\
		life' = life \land reportState' = reportState
	\end{schema}

	Exclusive backup publication uses the designated snapshot identity and the
	exact synced visible image.  It creates destination ownership and cleanup
	permission while retaining the source-stage obligation.

	% historical \begin{schema}{PublishBackupOwnership}
		\Delta Transaction \\
		TxnStaticFrame \\
		id? : Identity; delta? : Delta
		\where
		life = running \land cancelled = no \land phase = staging \land
		lockHeld = yes \land backupProgress = backupStaged \\
		id? \in backupIDs \cap sealedIDs \cap closeOK \\
		id? \in \dom syncedImages \cap \dom intendedImages \cap
		\dom snapshotImages \cap \dom visibleImages \\
		backupStage \in \dom visibleNS \cap \dom cleanupDue \land
		visibleNS~backupStage = file~id? \land cleanupDue~backupStage = id? \\
		visibleImages~id? = syncedImages~id? \land
		syncedImages~id? = intendedImages~id? \land
		intendedImages~id? = snapshotImages~id? \\
		imageState~(snapshotImages~id?) = complete \\
		backupPath \notin \dom visibleNS \\
		deltaRemove~delta? = \emptyset \land
		deltaPut~delta? = \{ backupPath \mapsto file~id? \} \\
		AppendPendingDelta \\
		visibleNS' = applyDelta~(visibleNS,delta?) \land durableNS' = durableNS \\
		visibleImages' = visibleImages \land syncedImages' = syncedImages \land
		intendedImages' = intendedImages \land snapshotImages' = snapshotImages \land
		dirtyChoices' = dirtyChoices \land stageIDs' = stageIDs \land
		sealedIDs' = sealedIDs \land closeOK' = closeOK \\
		published' = published \cup \{ backupPath \mapsto id? \} \\
		publicationImages' = publicationImages \oplus
		\{ id? \mapsto snapshotImages~id? \} \\
		owned' = owned \oplus \{ backupPath \mapsto id? \} \\
		cleanupDue' = cleanupDue \oplus \{ backupPath \mapsto id? \} \\
		cleanupPending' = cleanupPending \land retainedResults' = retainedResults \land
		cleanupNotes' = cleanupNotes \\
		phase' = publishingBackup \land committed' = committed \land
		configSyncAck' = configSyncAck \land cancelled' = cancelled \land
		lockHeld' = lockHeld \land backupProgress' = backupLinked \land
		life' = life \land reportState' = reportState
	\end{schema}

	% historical \begin{schema}{SyncBackupLink}
		\Delta Transaction \\
		TxnStaticFrame \\
		TxnImageFrame \\
		TxnOwnershipFrame
		\where
		life = running \land cancelled = no \land phase = publishingBackup \land
		lockHeld = yes \land backupProgress = backupLinked \\
		visibleNS' = visibleNS \land durableNS' = visibleNS \land
		pendingCount' = pendingZero \land
		pendingDeltaOne' = pendingDeltaOne \land
		pendingDeltaTwo' = pendingDeltaTwo \land
		pendingDeltaThree' = pendingDeltaThree \\
		phase' = publishingBackup \land committed' = committed \land
		configSyncAck' = configSyncAck \land cancelled' = cancelled \land
		lockHeld' = lockHeld \land backupProgress' = linkSynced \land
		life' = life \land reportState' = reportState
	\end{schema}

	% historical \begin{schema}{RemoveBackupStage}
		\Delta Transaction \\
		TxnStaticFrame \\
		TxnImageFrame \\
		id? : Identity; delta? : Delta
		\where
		life = running \land cancelled = no \land phase = publishingBackup \land
		lockHeld = yes \land backupProgress = linkSynced \land
		backupStage \in \dom cleanupDue \cap \dom visibleNS \\
		id? = cleanupDue~backupStage \land visibleNS~backupStage = file~id? \\
		backupStage \notin \dom retainedResults \\
		deltaRemove~delta? = \{ backupStage \} \land deltaPut~delta? = \emptyset \\
		AppendPendingDelta \\
		visibleNS' = applyDelta~(visibleNS,delta?) \land durableNS' = durableNS \\
		owned' = owned \land cleanupDue' = cleanupDue \land
		cleanupPending' = cleanupPending \oplus \{ backupStage \mapsto id? \} \\
		retainedResults' = retainedResults \land cleanupNotes' = cleanupNotes \\
		phase' = publishingBackup \land committed' = committed \land
		configSyncAck' = configSyncAck \land cancelled' = cancelled \land
		lockHeld' = lockHeld \land backupProgress' = stageRemoved \land
		life' = life \land reportState' = reportState
	\end{schema}

	% historical \begin{schema}{SyncBackupRemoval}
		\Delta Transaction \\
		TxnStaticFrame \\
		TxnImageFrame
		\where
		life = running \land cancelled = no \land phase = publishingBackup \land
		lockHeld = yes \land backupProgress = stageRemoved \land
		backupStage \notin \dom visibleNS \\
		visibleNS' = visibleNS \land durableNS' = visibleNS \land
		pendingCount' = pendingZero \land
		pendingDeltaOne' = pendingDeltaOne \land
		pendingDeltaTwo' = pendingDeltaTwo \land
		pendingDeltaThree' = pendingDeltaThree \\
		owned' = owned \land
		cleanupDue' = \{ backupStage \} \ndres cleanupDue \land
		cleanupPending' = \{ backupStage \} \ndres cleanupPending \land
		retainedResults' = retainedResults \land cleanupNotes' = cleanupNotes \\
		phase' = revalidating \land committed' = committed \land
		configSyncAck' = configSyncAck \land cancelled' = cancelled \land
		lockHeld' = lockHeld \land backupProgress' = backupReady \land
		life' = life \land reportState' = reportState
	\end{schema}

	The no-conflict path and the completed-backup path meet at an explicit
	revalidation phase.  A successful revalidation compares the accepted config
	existence/identity and complete image with the current namespace/image before
	exposing the publication phase.  The commit schemas require that milestone;
	neither cancellation nor abort can manufacture it.

	% historical \begin{schema}{BeginRevalidation}
		\Delta Transaction \\
		TxnStaticFrame \\
		TxnNamespaceFrame \\
		TxnImageFrame \\
		TxnOwnershipFrame
		\where
		life = running \land cancelled = no \land phase = staging \land
		lockHeld = yes \land conflicts = \emptyset \land
		backupProgress = backupNone \\
		candidateID \in sealedIDs \cap closeOK \cap \dom syncedImages \cap
		\dom intendedImages \\
		phase' = revalidating \land committed' = committed \land
		configSyncAck' = configSyncAck \land cancelled' = cancelled \land
		lockHeld' = lockHeld \land backupProgress' = backupProgress \land
		life' = life \land reportState' = reportState
	\end{schema}

	% historical \begin{schema}{RevalidationSuccess}
		\Delta Transaction \\
		TxnStaticFrame \\
		TxnNamespaceFrame \\
		TxnImageFrame \\
		TxnOwnershipFrame
		\where
		life = running \land cancelled = no \land phase = revalidating \land
		lockHeld = yes \\
		( expectedConfig = \emptyset \land configPath \notin \dom visibleNS ) \lor \\
		( \exists originalID : expectedConfig @ \\
		\t1 originalID \in \dom snapshotImages \cap \dom visibleImages \land \\
		\t1 \{ configPath \mapsto file~originalID \} \subseteq visibleNS \land
		visibleImages~originalID = snapshotImages~originalID ) \\
		phase' = publishingConfig \land committed' = committed \land
		configSyncAck' = configSyncAck \land cancelled' = cancelled \land
		lockHeld' = lockHeld \land backupProgress' = backupProgress \land
		life' = life \land reportState' = reportState
	\end{schema}

	Both commit operations require backup-ready exactly when conflicts exist.  They
	publish the complete candidate, retain the config destination, and atomically
	transfer the exact final backup from abort cleanup to retained result.  Rename
	records a source tombstone; hard-link publication leaves the source cleanup-due.

	% historical \begin{schema}{CommitExisting}
		\Delta Transaction \\
		TxnStaticFrame \\
		delta? : Delta; backupID? : Identity
		\where
		life = running \land cancelled = no \land phase = publishingConfig \land
		lockHeld = yes \land expectedConfig \neq \emptyset \land
		configPath \in \dom visibleNS \land
		visibleNS~configPath \neq file~candidateID \\
		( \exists originalID : expectedConfig @ \\
		\t1 originalID \in \dom snapshotImages \cap \dom visibleImages \land \\
		\t1 visibleNS~configPath = file~originalID \land
		visibleImages~originalID = snapshotImages~originalID ) \\
		configStage \in \dom visibleNS \cap \dom cleanupDue \land
		visibleNS~configStage = file~candidateID \land
		cleanupDue~configStage = candidateID \\
		candidateID \in sealedIDs \cap closeOK \cap \dom syncedImages \cap
		\dom intendedImages \cap \dom visibleImages \\
		visibleImages~candidateID = syncedImages~candidateID \land
		syncedImages~candidateID = intendedImages~candidateID \land
		imageState~(intendedImages~candidateID) = complete \\
		( conflicts = \emptyset \land backupProgress = backupNone \land
		backupPath \notin \dom owned ) \lor
		( conflicts \neq \emptyset \land backupProgress = backupReady \land
		backupPath \in \dom cleanupDue \cap \dom visibleNS \land
		backupID? = cleanupDue~backupPath \land backupID? \in backupIDs \land
		visibleNS~backupPath = file~backupID? ) \\
		deltaRemove~delta? = \{ configStage \} \land
		deltaPut~delta? = \{ configPath \mapsto file~candidateID \} \\
		AppendPendingDelta \\
		visibleNS' = applyDelta~(visibleNS,delta?) \land durableNS' = durableNS \\
		visibleImages' = visibleImages \land syncedImages' = syncedImages \land
		intendedImages' = intendedImages \land snapshotImages' = snapshotImages \land
		dirtyChoices' = dirtyChoices \land stageIDs' = stageIDs \land
		sealedIDs' = sealedIDs \land closeOK' = closeOK \\
		published' = published \cup \{ configPath \mapsto candidateID \} \\
		publicationImages' = publicationImages \oplus
		\{ candidateID \mapsto intendedImages~candidateID \} \\
		owned' = owned \oplus \{ configPath \mapsto candidateID \} \\
		cleanupDue' = \{ backupPath \} \ndres cleanupDue \\
		cleanupPending' = cleanupPending \oplus \{ configStage \mapsto candidateID \} \\
		retainedResults' = retainedResults \oplus
		\{ configPath \mapsto candidateID \} \oplus
		( \{ backupPath \} \dres owned ) \\
		cleanupNotes' = cleanupNotes \\
		phase' = confirming \land committed' = yes \land configSyncAck' = no \land
		cancelled' = cancelled \land lockHeld' = lockHeld \land
		backupProgress' = backupProgress \land life' = life \land
		reportState' = reportState
	\end{schema}

	Cancellation is observed at a boundary before the next live operation.  Before
	commit it enters $aborting$ immediately; after commit it records cancellation
	without blocking mandatory durability confirmation and identity-safe cleanup.
	No forward staging, backup, revalidation, or commit transition accepts either
	$cancelled=yes$ or $phase=aborting$.

	% historical \begin{schema}{ObserveCancellation}
		\Delta Transaction \\
		TxnStaticFrame \\
		TxnNamespaceFrame \\
		TxnImageFrame \\
		TxnOwnershipFrame
		\where
		life = running \land cancelled = no \land phase \neq finished \land
		reportState = reportPending \\
		( committed = no \implies phase' = aborting ) \\
		( committed = yes \implies phase' = phase ) \\
		committed' = committed \land configSyncAck' = configSyncAck \land
		cancelled' = yes \land lockHeld' = lockHeld \land
		backupProgress' = backupProgress \land life' = life \land
		reportState' = reportState
	\end{schema}

	% historical \begin{schema}{CommitAbsent}
		\Delta Transaction \\
		TxnStaticFrame \\
		delta? : Delta; backupID? : Identity
		\where
		life = running \land cancelled = no \land phase = publishingConfig \land
		lockHeld = yes \land expectedConfig = \emptyset \land
		configPath \notin \dom visibleNS \\
		configStage \in \dom visibleNS \cap \dom cleanupDue \land
		visibleNS~configStage = file~candidateID \land
		cleanupDue~configStage = candidateID \\
		candidateID \in sealedIDs \cap closeOK \cap \dom syncedImages \cap
		\dom intendedImages \cap \dom visibleImages \\
		visibleImages~candidateID = syncedImages~candidateID \land
		syncedImages~candidateID = intendedImages~candidateID \land
		imageState~(intendedImages~candidateID) = complete \\
		( conflicts = \emptyset \land backupProgress = backupNone \land
		backupPath \notin \dom owned ) \lor
		( conflicts \neq \emptyset \land backupProgress = backupReady \land
		backupPath \in \dom cleanupDue \cap \dom visibleNS \land
		backupID? = cleanupDue~backupPath \land backupID? \in backupIDs \land
		visibleNS~backupPath = file~backupID? ) \\
		deltaRemove~delta? = \emptyset \land
		deltaPut~delta? = \{ configPath \mapsto file~candidateID \} \\
		AppendPendingDelta \\
		visibleNS' = applyDelta~(visibleNS,delta?) \land durableNS' = durableNS \\
		visibleImages' = visibleImages \land syncedImages' = syncedImages \land
		intendedImages' = intendedImages \land snapshotImages' = snapshotImages \land
		dirtyChoices' = dirtyChoices \land stageIDs' = stageIDs \land
		sealedIDs' = sealedIDs \land closeOK' = closeOK \\
		published' = published \cup \{ configPath \mapsto candidateID \} \\
		publicationImages' = publicationImages \oplus
		\{ candidateID \mapsto intendedImages~candidateID \} \\
		owned' = owned \oplus \{ configPath \mapsto candidateID \} \\
		cleanupDue' = \{ backupPath \} \ndres cleanupDue \\
		cleanupPending' = cleanupPending \\
		retainedResults' = retainedResults \oplus
		\{ configPath \mapsto candidateID \} \oplus
		( \{ backupPath \} \dres owned ) \\
		cleanupNotes' = cleanupNotes \\
		phase' = confirming \land committed' = yes \land configSyncAck' = no \land
		cancelled' = cancelled \land lockHeld' = lockHeld \land
		backupProgress' = backupProgress \land life' = life \land
		reportState' = reportState
	\end{schema}

	% historical \begin{schema}{UnlinkOwned}
		\Delta Transaction \\
		TxnStaticFrame \\
		TxnImageFrame \\
		TxnControlFrame \\
		p? : Path; id? : Identity; delta? : Delta
		\where
		life = running \land phase \in \{ aborting, confirming, cleaning \} \land
		lockHeld = yes \\
		p? \in \dom cleanupDue \cap \dom visibleNS \\
		p? \notin \dom retainedResults \land id? = cleanupDue~p? \land
		visibleNS~p? = file~id? \\
		deltaRemove~delta? = \{ p? \} \land deltaPut~delta? = \emptyset \\
		AppendPendingDelta \\
		visibleNS' = applyDelta~(visibleNS,delta?) \land durableNS' = durableNS \\
		owned' = owned \land cleanupDue' = cleanupDue \land
		cleanupPending' = cleanupPending \oplus \{ p? \mapsto id? \} \\
		retainedResults' = retainedResults \land cleanupNotes' = cleanupNotes
	\end{schema}

	% historical \begin{schema}{CleanupUnlinkFailed}
		\Delta Transaction \\
		TxnStaticFrame \\
		TxnNamespaceFrame \\
		TxnImageFrame \\
		TxnControlFrame \\
		p? : Path
		\where
		life = running \land phase \in \{ aborting, confirming, cleaning \} \land
		lockHeld = yes \\
		p? \in \dom cleanupDue \setminus \dom retainedResults \\
		owned' = owned \land cleanupDue' = cleanupDue \land
		cleanupPending' = cleanupPending \land retainedResults' = retainedResults \land
		cleanupNotes' = cleanupNotes \oplus \{ p? \mapsto unlinkFailed \}
	\end{schema}

	% historical \begin{schema}{CleanupReplacementFound}
		\Delta Transaction \\
		TxnStaticFrame \\
		TxnNamespaceFrame \\
		TxnImageFrame \\
		TxnControlFrame \\
		p? : Path; replacement? : Identity
		\where
		life = running \land phase \in \{ aborting, confirming, cleaning \} \land
		lockHeld = yes \\
		p? \in \dom cleanupDue \cap \dom visibleNS \\
		visibleNS~p? = file~replacement? \land replacement? \neq cleanupDue~p? \\
		owned' = owned \land cleanupDue' = cleanupDue \land
		cleanupPending' = cleanupPending \land retainedResults' = retainedResults \land
		cleanupNotes' = cleanupNotes \oplus \{ p? \mapsto identityChanged \}
	\end{schema}

	% historical \begin{schema}{CleanupInspectionFailed}
		\Delta Transaction \\
		TxnStaticFrame \\
		TxnNamespaceFrame \\
		TxnImageFrame \\
		TxnControlFrame \\
		p? : Path
		\where
		life = running \land phase \in \{ aborting, confirming, cleaning \} \land
		lockHeld = yes \\
		p? \in \dom cleanupDue \\
		owned' = owned \land cleanupDue' = cleanupDue \land
		cleanupPending' = cleanupPending \land retainedResults' = retainedResults \land
		cleanupNotes' = cleanupNotes \oplus \{ p? \mapsto inspectFailed \}
	\end{schema}

	% historical \begin{schema}{DirectorySyncSuccess}
		\Delta Transaction \\
		TxnStaticFrame \\
		TxnImageFrame
		\where
		life = running \land phase \in \{ aborting, confirming, cleaning \} \land
		lockHeld = yes \\
		visibleNS' = visibleNS \land durableNS' = visibleNS \land
		pendingCount' = pendingZero \land
		pendingDeltaOne' = pendingDeltaOne \land
		pendingDeltaTwo' = pendingDeltaTwo \land
		pendingDeltaThree' = pendingDeltaThree \\
		owned' = owned \\
		cleanupDue' = \{ p : \dom cleanupDue |
		\lnot ( \{ p \mapsto file~(cleanupDue~p) \} \subseteq visibleNS ) @ p \}
		\ndres cleanupDue \\
		cleanupPending' = \{ p : \dom cleanupPending |
		\lnot ( \{ p \mapsto file~(cleanupDue~p) \} \subseteq visibleNS ) @ p \}
		\ndres cleanupPending \\
		retainedResults' = retainedResults \land cleanupNotes' = cleanupNotes \\
		phase' = phase \land committed' = committed \land
		( committed = yes \land
		\{ configPath \mapsto file~candidateID \} \subseteq visibleNS
		\implies configSyncAck' = yes ) \\
		( committed = yes \land
		\lnot ( \{ configPath \mapsto file~candidateID \} \subseteq visibleNS ) \implies
		configSyncAck' = configSyncAck ) \\
		( committed = no \implies configSyncAck' = configSyncAck ) \\
		cancelled' = cancelled \land lockHeld' = lockHeld \land
		backupProgress' = backupProgress \land life' = life \land
		reportState' = reportState
	\end{schema}

	% historical \begin{schema}{DirectorySyncFailure}
		\Xi Transaction
		\where
		life = running \land phase \in \{ aborting, confirming, cleaning \} \land
		lockHeld = yes
	\end{schema}

	% historical \begin{theorem}{ConfirmedReplacementDischargesOwnedIdentity}
		\forall DirectorySyncSuccess; p : \dom cleanupDue \cap \dom cleanupNotes @ \\
		\t1 cleanupNotes~p = identityChanged \land
		\lnot ( \{ p \mapsto file~(cleanupDue~p) \} \subseteq visibleNS ) \\
		\t2 \implies p \notin \dom cleanupDue' \land
		cleanupNotes'~p = identityChanged
	\end{theorem}

	% historical \begin{schema}{AbortTransaction}
		\Delta Transaction \\
		TxnStaticFrame \\
		TxnNamespaceFrame \\
		TxnImageFrame \\
		TxnOwnershipFrame
		\where
		committed = no \land life = running \land
		phase \notin \{ aborting, finished \} \\
		phase' = aborting \land committed' = no \land
		configSyncAck' = configSyncAck \land cancelled' = cancelled \land
		lockHeld' = lockHeld \land backupProgress' = backupProgress \land
		life' = life \land reportState' = reportState
	\end{schema}

	\subsection{Finite paired crash restoration}

	The crash relation chooses each active delta or skips it, in fixed slot order,
	then chooses exactly one actually reached finite image for each referenced
	identity.  This includes the empty and full choices without a recursive
	powerset-valued function.  Namespace and images are one pair.

	% historical \begin{schema}{CrashPair}
		Transaction \\
		crashNS : Namespace \\
		crashImages : Identity \pfun FileImage \\
		takeOne, takeTwo, takeThree : Bool \\
		afterOne, afterTwo : Namespace
		\where
		( pendingCount = pendingZero \implies
		takeOne = no \land takeTwo = no \land takeThree = no ) \\
		( pendingCount = pendingOne \implies takeTwo = no \land takeThree = no ) \\
		( pendingCount = pendingTwo \implies takeThree = no ) \\
		( takeOne = yes \implies pendingCount \neq pendingZero \land
		afterOne = applyDelta~(durableNS,pendingDeltaOne) ) \\
		( takeOne = no \implies afterOne = durableNS ) \\
		( takeTwo = yes \implies pendingCount \in \{ pendingTwo, pendingThree \} \land
		afterTwo = applyDelta~(afterOne,pendingDeltaTwo) ) \\
		( takeTwo = no \implies afterTwo = afterOne ) \\
		( takeThree = yes \implies pendingCount = pendingThree \land
		crashNS = applyDelta~(afterTwo,pendingDeltaThree) ) \\
		( takeThree = no \implies crashNS = afterTwo ) \\
		\dom crashImages = refs~crashNS \\
		( \forall id : refs~crashNS @ \\
		\t1 ( id \in \dom syncedImages \land crashImages~id = syncedImages~id ) \lor \\
		\t1 ( id \in \dom dirtyChoices \land crashImages~id \in dirtyChoices~id ) ) \\
		( \forall p : Path; id : Identity | (p,id) \in published \land \\
		\t1 p \in \dom crashNS \land crashNS~p = file~id @ \\
		\t2 crashImages~id = publicationImages~id )
	\end{schema}

	% historical \begin{schema}{PowerLoss}
		\Delta Transaction \\
		CrashPair
		\where
		life = running \land phase \neq finished \\
		visibleNS' = crashNS \\
		durableNS' = crashNS \\
		pendingCount' = pendingZero \\
		pendingDeltaOne' = pendingDeltaOne \\
		pendingDeltaTwo' = pendingDeltaTwo \\
		pendingDeltaThree' = pendingDeltaThree \\
		visibleImages' = crashImages \\
		syncedImages' = crashImages \\
		dirtyChoices' = \emptyset \\
		intendedImages' = intendedImages \\
		snapshotImages' = snapshotImages \\
		stageIDs' = \emptyset \\
		sealedIDs' = \emptyset \\
		closeOK' = \emptyset \\
		owned' = \emptyset \\
		cleanupDue' = \emptyset \\
		cleanupPending' = \emptyset \\
		retainedResults' = \emptyset \\
		cleanupNotes' = \emptyset \\
		published' = \{ p : \dom crashNS; id : Identity |
		(p,id) \in published \land crashNS~p = file~id @ p \mapsto id \} \\
		publicationImages' = publicationImages \\
		configPath' = configPath \\
		backupPath' = backupPath \\
		configStage' = configStage \\
		backupStage' = backupStage \\
		candidateID' = candidateID \\
		backupIDs' = backupIDs \\
		expectedConfig' = expectedConfig \\
		conflicts' = conflicts \\
		phase' = finished \\
		committed' = committed \\
		configSyncAck' = configSyncAck \\
		cancelled' = cancelled \\
		lockHeld' = no \\
		backupProgress' = backupProgress \\
		life' = ended \\
		reportState' = processDead
	\end{schema}

	% historical \begin{schema}{ProcessDeath}
		\Delta Transaction \\
		TxnStaticFrame \\
		TxnNamespaceFrame \\
		TxnImageFrame
		\where
		life = running \land phase \neq finished \\
		owned' = \emptyset \land cleanupDue' = \emptyset \land
		cleanupPending' = \emptyset \land retainedResults' = \emptyset \land
		cleanupNotes' = \emptyset \\
		phase' = finished \land committed' = committed \land
		configSyncAck' = configSyncAck \land cancelled' = cancelled \land
		lockHeld' = no \land backupProgress' = backupProgress \land
		life' = ended \land reportState' = processDead
	\end{schema}

	The complete transition additionally clears the active pending count, dirty
	choices, cleanup notes, stage/seal/close sets and lock ownership.  Assigning restored images as a
	new settled baseline does not turn their preserved $partial$ marker into
	$complete$.  A restarted process receives no dead invocation's cleanup rights.
	Process death without power loss ends execution and releases the lock while
	leaving the visible filesystem unchanged.

	% historical \begin{theorem}{PublishedCrashImageExact}
		\forall CrashPair; p : Path; id : Identity @ \\
		\t1 (p,id) \in published \land p \in \dom crashNS \land crashNS~p = file~id \\
		\t2 \implies crashImages~id = publicationImages~id
	\end{theorem}

	% historical \begin{theorem}{CrashImageCoverage}
		\forall CrashPair @ \dom crashImages = refs~crashNS
	\end{theorem}

	% historical \begin{theorem}{PublishedImagesAreComplete}
		\forall Transaction; id : Identity | id \in \dom publicationImages @
		imageState~(publicationImages~id) = complete
	\end{theorem}

	% historical \begin{theorem}{PublishedBackupIsExactSnapshot}
		\forall Transaction; id : backupIDs \cap \dom publicationImages @
		publicationImages~id = snapshotImages~id
	\end{theorem}

	% historical \begin{theorem}{DirtyChoicesContainNoFutureBytes}
		\forall Transaction; id : \dom dirtyChoices; i : dirtyChoices~id @ \\
		\t1 imageBytes~i \mapsto imageBytes~(visibleImages~id) \in bytePrefix \land \\
		\t1 imageBytes~(visibleImages~id) \mapsto
		imageBytes~(intendedImages~id) \in bytePrefix
	\end{theorem}

	For cleanup reporting, owned-present means every $CrashPair$ contains that
	identity at the path; owned-absent means every pair omits it there; otherwise
	presence is uncertain.  A visible unlink by itself proves neither durable
	absence nor config durability.

	The required bounded traces are: cancellation before publication entering abort;
	cancellation after publication followed by mandatory sync and cleanup; successful
	revalidation before both commit modes; abort after backup-ready; backup
	replacement followed by sync that discharges only the owned identity while
	retaining its note; unlink with failed sync and crash resurrection; rename-source
	tombstone; absent-config hard link with only the stage cleaned; conflicting
	commit transfer which later cancellation/sync/cleanup/report failure cannot
	reverse; fresh empty, partial, and complete-unsynced stage crash images with no
	future bytes; directory sync without file sync; final backup surviving without
	its stage; complete final config after rename/link; exact snapshot backup
	bytes/mode; external dirty original choices; shared hard-link image identity;
	and process death/restart without a successor write, commit, or inherited
	authority.  Result witnesses include a retained $dConfigSync$ event followed by
	confirmed acknowledgement and a failed-after-commit result, plus both early
	null-URL and post-local derived-URL cancellation.

	Mutation checking must obtain counterexamples after removing file-sync before
	publication, destination ownership/cleanup insertion, rename-source tombstone,
	backup cleanup-to-retained transfer, cleanup directory-sync acknowledgement,
	replacement-identity discharge, successful revalidation/cancellation guards,
	reached-prefix image constraints, published-image restoration equality,
	backup-before-commit, HTTP authorization, or operational-event precedence.  Such
	mutants are deliberately not alternative legal schemas; they are checker
	obligations against the predicates above.

\end{verbatim}

\subsection{Increment 1 executable transaction library}

The raw-routing $Init$ above remains quiescent: it initializes only
$RawConfig$ and does not start a filesystem invocation.  The declarations and
safety predicate below are an independent library.  They mention no
$RawConfig$ component, and no transaction operation includes $TxnSafety'$ as a
postcondition.  Only a generated finite scenario introduces the ProZ names
$State$, $Init$, and $Invariant$.

% POINTS-CHECK BEGIN types.txn type
\begin{zed}
	EntryRef ::= file \ldata Identity \rdata | symlinkRef | directoryRef | specialRef \\
	Namespace == Path \pfun EntryRef \\
	CleanupObservation ::= noCleanupNote | unlinkFailed | identityChanged | inspectFailed \\
	BackupProgress ::= backupNone | backupStaged | backupLinked | linkSynced | \\
	\t1 stageRemoved | backupReady \\
	PendingCount ::= pendingZero | pendingOne | pendingTwo | pendingThree | \\
	\t1 pendingFour | pendingFive | pendingSix \\
	TransactionPhase ::= validating | staging | publishingBackup | revalidating | \\
	\t1 publishingConfig | confirming | cleaning | aborting | sweeping | \\
	\t1 cleanupSync | releasing | reporting | finished \\
	Life ::= running | ended \\
	ReportState ::= reportPending | reportDelivered | reportFailed | processDead \\
	LinkCount ::= oneLink | multipleLinks \\
	AccessTag ::= supportedAccessA | supportedAccessB | unsupportedAccess
\end{zed}
% POINTS-CHECK END types.txn

% POINTS-CHECK BEGIN definitions.filesystem definition
\begin{axdef}
	deltaRemove : Delta \fun \power Path \\
	deltaPut : Delta \fun Namespace \\
	emptyBytes : ByteString \\
	bytePrefix : ByteString \rel ByteString \\
	imageBytes : FileImage \fun ByteString \\
	imageMode : FileImage \fun Mode \\
	imageOwner : FileImage \fun Owner \\
	imageState : FileImage \fun ContentState
	\where
	\forall b : ByteString @ emptyBytes \mapsto b \in bytePrefix \\
	\forall b : ByteString @ b \mapsto b \in bytePrefix \\
	\forall a, b : ByteString @ a \mapsto b \in bytePrefix \land
	b \mapsto a \in bytePrefix \implies a = b \\
	\forall a, b, c : ByteString @ a \mapsto b \in bytePrefix \land
	b \mapsto c \in bytePrefix \implies a \mapsto c \in bytePrefix \\
	\forall b : ByteString; m : Mode; o : Owner; c : ContentState @ \\
	\t1 \exists_1 i : FileImage @ imageBytes~i = b \land imageMode~i = m \land \\
	\t2 imageOwner~i = o \land imageState~i = c \\
	\forall i, j : FileImage @ i = j \iff \\
	\t1 imageBytes~i = imageBytes~j \land imageMode~i = imageMode~j \land \\
	\t1 imageOwner~i = imageOwner~j \land imageState~i = imageState~j
\end{axdef}

\begin{zed}
	applyDelta == ( \lambda d : Delta @ ( \lambda n : Namespace @ \\
	\t1 (deltaRemove~d \ndres n) \oplus deltaPut~d ) ) \\
	refs == ( \lambda n : Namespace @ \\
	\t1 \{ id : Identity | \exists p : \dom n @ n~p = file~id \} )
\end{zed}
% POINTS-CHECK END definitions.filesystem

% POINTS-CHECK BEGIN state.txnvars state
\begin{schema}{TxnVars}
	visibleNS, durableNS : Namespace \\
	pendingCount : PendingCount \\
	pendingDeltaOne, pendingDeltaTwo, pendingDeltaThree : Delta \\
	pendingDeltaFour, pendingDeltaFive, pendingDeltaSix : Delta \\
	visibleImages, syncedImages, intendedImages, snapshotImages :
	Identity \pfun FileImage \\
	dirtyChoices : Identity \pfun \power FileImage \\
	stageIDs, sealedIDs, closeOK : \power Identity \\
	published : \power ( Path \cross Identity ) \\
	publicationImages : Identity \pfun FileImage \\
	owned, cleanupDue, cleanupPending, retainedResults : Path \pfun Identity \\
	cleanupNotes : Path \pfun CleanupObservation \\
	configPath, backupPath, configStage, backupStage : Path \\
	candidateID : Identity \\
	backupIDs, expectedConfig : \power Identity \\
	conflicts : \power JSONPointer \\
	phase : TransactionPhase \\
	committed, configSyncAck, cancelled, lockHeld : Bool \\
	backupProgress : BackupProgress \\
	life : Life \\
	reportState : ReportState \\
	revalidationPassed, confirmationAttempted : Bool \\
	cleanupWork : Path \pfun Identity \\
	cleanupTried : \power Path \\
	cleanupSyncAttempted, releaseAttempted : Bool \\
	events : \power RouteEvent \\
	inputAccepted, writeStarted, noOp, namespaceMutated : Bool \\
	snapshotLinks, currentLinks : Identity \pfun LinkCount \\
	snapshotAccess, currentAccess : Identity \pfun AccessTag \\
	ownedBackupPublication : Path \pfun Identity \\
	backupAbsenceConfirmed : Bool \\
	authorizationPassed : Bool \\
	authorizedTransport : \power Transport \\
	insecureConsent : Bool
\end{schema}
% POINTS-CHECK END state.txnvars

The following predicate is checked independently as the generated machine
invariant.  Active-delta coverage is explicit for all six slots; no recursive
sequence or powerset-of-sequences construction is used.

% POINTS-CHECK BEGIN predicate.txnsafety predicate
\begin{schema}{TxnSafety}
	TxnVars
	\where
	\# \{ configPath, backupPath, configStage, backupStage \} = 4 \\
	\# expectedConfig \leq 1 \\
	candidateID \notin backupIDs \\
	cleanupDue \subseteq owned \\
	cleanupPending \subseteq cleanupDue \\
	retainedResults \subseteq owned \\
	\dom cleanupDue \cap \dom retainedResults = \emptyset \\
	\dom syncedImages \cap \dom dirtyChoices = \emptyset \\
	( \forall choices : \ran dirtyChoices @ choices \neq \emptyset ) \\
	( \forall id : \dom intendedImages @
	imageState~(intendedImages~id) = complete ) \\
	( \forall id : \dom dirtyChoices @ \\
	\t1 id \in \dom visibleImages \cap \dom intendedImages \land \\
	\t1 imageBytes~(visibleImages~id) \mapsto
	imageBytes~(intendedImages~id) \in bytePrefix \land \\
	\t1 dirtyChoices~id = \{ i : FileImage | \\
	\t2 imageBytes~i \mapsto imageBytes~(visibleImages~id) \in bytePrefix \land \\
	\t2 imageMode~i = imageMode~(intendedImages~id) \land \\
	\t2 imageOwner~i = imageOwner~(intendedImages~id) \land \\
	\t2 ( imageState~i = complete \iff
	imageBytes~i = imageBytes~(intendedImages~id) ) @ i \} ) \\
	refs~visibleNS \cup refs~durableNS \subseteq
	\dom syncedImages \cup \dom dirtyChoices \\
	( pendingCount \in \{ pendingOne, pendingTwo, pendingThree, pendingFour,
	pendingFive, pendingSix \} \implies
	refs~(deltaPut~pendingDeltaOne) \subseteq
	\dom syncedImages \cup \dom dirtyChoices ) \\
	( pendingCount \in \{ pendingTwo, pendingThree, pendingFour, pendingFive,
	pendingSix \} \implies refs~(deltaPut~pendingDeltaTwo) \subseteq
	\dom syncedImages \cup \dom dirtyChoices ) \\
	( pendingCount \in \{ pendingThree, pendingFour, pendingFive, pendingSix \}
	\implies refs~(deltaPut~pendingDeltaThree) \subseteq
	\dom syncedImages \cup \dom dirtyChoices ) \\
	( pendingCount \in \{ pendingFour, pendingFive, pendingSix \} \implies
	refs~(deltaPut~pendingDeltaFour) \subseteq
	\dom syncedImages \cup \dom dirtyChoices ) \\
	( pendingCount \in \{ pendingFive, pendingSix \} \implies
	refs~(deltaPut~pendingDeltaFive) \subseteq
	\dom syncedImages \cup \dom dirtyChoices ) \\
	( pendingCount = pendingSix \implies
	refs~(deltaPut~pendingDeltaSix) \subseteq
	\dom syncedImages \cup \dom dirtyChoices ) \\
	( \forall p : Path; id : Identity | (p,id) \in published @ \\
	\t1 id \in \dom publicationImages \land
	imageState~(publicationImages~id) = complete ) \\
	( committed = yes \implies revalidationPassed = yes ) \\
	( configSyncAck = yes \implies committed = yes ) \\
	( life = ended \implies phase = finished \land lockHeld = no ) \\
	( \dom retainedResults \cap \dom cleanupDue = \emptyset )
\end{schema}
% POINTS-CHECK END predicate.txnsafety

% POINTS-CHECK BEGIN predicate.snapshot predicate
\begin{schema}{SnapshotMatches}
	TxnVars
	\where
	( expectedConfig = \emptyset \land configPath \notin \dom visibleNS ) \lor \\
	( \exists originalID : expectedConfig @ \\
	\t1 \{ configPath \mapsto file~originalID \} \subseteq visibleNS \land \\
	\t1 originalID \in \dom visibleImages \cap \dom snapshotImages \land \\
	\t1 visibleImages~originalID = snapshotImages~originalID \land \\
	\t1 originalID \in \dom currentLinks \cap \dom snapshotLinks \land \\
	\t1 currentLinks~originalID = snapshotLinks~originalID \land \\
	\t1 originalID \in \dom currentAccess \cap \dom snapshotAccess \land \\
	\t1 currentAccess~originalID = snapshotAccess~originalID )
\end{schema}
% POINTS-CHECK END predicate.snapshot

% POINTS-CHECK BEGIN predicate.backup predicate
\begin{schema}{BackupBarrier}
	TxnVars
	\where
	( conflicts = \emptyset \land backupProgress = backupNone \land
	backupPath \notin \dom owned ) \lor \\
	( conflicts \neq \emptyset \land backupProgress = backupReady \land \\
	\t1 ( \exists id : backupIDs @ \\
	\t2 \{ backupPath \mapsto id \} \subseteq ownedBackupPublication \land \\
	\t2 \{ backupPath \mapsto id \} \subseteq cleanupDue \land \\
	\t2 \{ backupPath \mapsto file~id \} \subseteq visibleNS \cap durableNS ) \land \\
	\t1 backupStage \notin \dom visibleNS \cup \dom durableNS )
\end{schema}

\begin{schema}{RevalidationBarrier}
	TxnVars
	\where
	phase = publishingConfig \land revalidationPassed = yes
\end{schema}
% POINTS-CHECK END predicate.backup

% POINTS-CHECK BEGIN predicate.publication predicate
\begin{schema}{WrittenComplete}
	TxnVars \\
	id? : Identity
	\where
	id? \in \dom visibleImages \cap \dom intendedImages \\
	visibleImages~id? = intendedImages~id? \\
	imageState~(intendedImages~id?) = complete
\end{schema}

\begin{schema}{FileSynced}
	TxnVars \\
	id? : Identity
	\where
	id? \in sealedIDs \cap \dom syncedImages \cap \dom intendedImages \\
	syncedImages~id? = intendedImages~id?
\end{schema}

\begin{schema}{PublicationQualified}
	TxnVars \\
	id? : Identity
	\where
	WrittenComplete \\
	% POINTS-MUTATION BEGIN m2.publication.file-sync
	FileSynced \\
	% POINTS-MUTATION END m2.publication.file-sync
	id? \in closeOK
\end{schema}
% POINTS-CHECK END predicate.publication

The common accepted-write initializer starts after authorization, safe locked
observation, successful normalization and serialization, and the non-no-op
decision.  Those facts are the $AcceptedWriteFacts$ of the refinement; the
scenario suffix supplies $InitFromAcceptedWrite$ by assigning every field not
fixed here.  In particular it binds exact candidate bytes, exact conflicts,
path roles, transport authorization, and the accepted existence observation.

% POINTS-CHECK BEGIN init.txn init
\begin{schema}{TxnInit}
	TxnVars'
	\where
	phase' = staging \land life' = running \land reportState' = reportPending \\
	lockHeld' = yes \land inputAccepted' = yes \land writeStarted' = yes \land
	noOp' = no \\
	committed' = no \land configSyncAck' = no \land cancelled' = no \\
	revalidationPassed' = no \land confirmationAttempted' = no \\
	cleanupSyncAttempted' = no \land releaseAttempted' = no \\
	namespaceMutated' = no \land backupAbsenceConfirmed' = no \land
	backupProgress' = backupNone \\
	events' = \emptyset \\
	stageIDs' = \emptyset \land sealedIDs' = \emptyset \land closeOK' = \emptyset \\
	owned' = \emptyset \land cleanupDue' = \emptyset \land
	cleanupPending' = \emptyset \land retainedResults' = \emptyset \\
	cleanupNotes' = \emptyset \land cleanupWork' = \emptyset \land
	cleanupTried' = \emptyset \\
	published' = \emptyset \land publicationImages' = \emptyset \land
	ownedBackupPublication' = \emptyset \\
	pendingCount' = pendingZero \\
	pendingDeltaOne' = pendingDeltaTwo' = pendingDeltaThree' \\
	pendingDeltaFour' = pendingDeltaFive' = pendingDeltaSix'
\end{schema}
% POINTS-CHECK END init.txn

The authoritative $TxnInit$ schema remains in the reviewed source inventory.
The Scenario~A generator hashes its exact predicate body and inlines those
predicates into the sole generated initializer; it does not project the source
schema or export an initializer helper to ProZ.

The final initializer equality is completed by Scenario~A with all six slots
equal to its distinguished idle delta.  No field of the generated $Init$ is
unconstrained.

% POINTS-CHECK BEGIN effect.frames definition
\begin{schema}{TxnStaticFrame}
	\Delta TxnVars
	\where
	configPath' = configPath \land backupPath' = backupPath \land
	configStage' = configStage \land backupStage' = backupStage \\
	candidateID' = candidateID \land backupIDs' = backupIDs \land
	expectedConfig' = expectedConfig \land conflicts' = conflicts \\
	intendedImages' = intendedImages \land snapshotImages' = snapshotImages \\
	snapshotLinks' = snapshotLinks \land currentLinks' = currentLinks \land
	snapshotAccess' = snapshotAccess \land currentAccess' = currentAccess \\
	inputAccepted' = inputAccepted \land writeStarted' = writeStarted \land
	noOp' = noOp \\
	authorizationPassed' = authorizationPassed \land
	authorizedTransport' = authorizedTransport \land
	insecureConsent' = insecureConsent
\end{schema}

\begin{schema}{TxnNamespaceFrame}
	\Delta TxnVars
	\where
	visibleNS' = visibleNS \land durableNS' = durableNS \land
	pendingCount' = pendingCount \\
	pendingDeltaOne' = pendingDeltaOne \land pendingDeltaTwo' = pendingDeltaTwo \land
	pendingDeltaThree' = pendingDeltaThree \land
	pendingDeltaFour' = pendingDeltaFour \land
	pendingDeltaFive' = pendingDeltaFive \land pendingDeltaSix' = pendingDeltaSix \\
	namespaceMutated' = namespaceMutated
\end{schema}

\begin{schema}{TxnImageFrame}
	\Delta TxnVars
	\where
	visibleImages' = visibleImages \land syncedImages' = syncedImages \land
	dirtyChoices' = dirtyChoices
\end{schema}

\begin{schema}{TxnStageFrame}
	\Delta TxnVars
	\where
	stageIDs' = stageIDs \land sealedIDs' = sealedIDs \land closeOK' = closeOK
\end{schema}

\begin{schema}{TxnPublicationFrame}
	\Delta TxnVars
	\where
	published' = published \land publicationImages' = publicationImages
\end{schema}

\begin{schema}{TxnOwnershipFrame}
	\Delta TxnVars
	\where
	owned' = owned \land cleanupDue' = cleanupDue \land
	cleanupPending' = cleanupPending \land retainedResults' = retainedResults \land
	cleanupNotes' = cleanupNotes
\end{schema}

\begin{schema}{TxnControlFrame}
	\Delta TxnVars
	\where
	phase' = phase \land committed' = committed \land
	configSyncAck' = configSyncAck \land cancelled' = cancelled \land
	lockHeld' = lockHeld \land backupProgress' = backupProgress \land
	life' = life \land reportState' = reportState
\end{schema}

\begin{schema}{TxnAttemptFrame}
	\Delta TxnVars
	\where
	revalidationPassed' = revalidationPassed \land
	confirmationAttempted' = confirmationAttempted \land
	cleanupWork' = cleanupWork \land cleanupTried' = cleanupTried \land
	cleanupSyncAttempted' = cleanupSyncAttempted \land
	releaseAttempted' = releaseAttempted
\end{schema}

\begin{schema}{TxnEventFrame}
	\Delta TxnVars
	\where events' = events
\end{schema}

\begin{schema}{TxnBackupHistoryFrame}
	\Delta TxnVars
	\where
	ownedBackupPublication' = ownedBackupPublication \land
	backupAbsenceConfirmed' = backupAbsenceConfirmed
\end{schema}
% POINTS-CHECK END effect.frames

% POINTS-CHECK BEGIN effect.append definition
\begin{schema}{AppendPendingDelta}
	pendingCount, pendingCount' : PendingCount \\
	pendingDeltaOne, pendingDeltaTwo, pendingDeltaThree : Delta \\
	pendingDeltaFour, pendingDeltaFive, pendingDeltaSix : Delta \\
	pendingDeltaOne', pendingDeltaTwo', pendingDeltaThree' : Delta \\
	pendingDeltaFour', pendingDeltaFive', pendingDeltaSix' : Delta \\
	delta? : Delta
	\where
	( pendingCount = pendingZero \land pendingCount' = pendingOne \land
	pendingDeltaOne' = delta? \land pendingDeltaTwo' = pendingDeltaTwo \land
	pendingDeltaThree' = pendingDeltaThree \land pendingDeltaFour' = pendingDeltaFour \land
	pendingDeltaFive' = pendingDeltaFive \land pendingDeltaSix' = pendingDeltaSix ) \lor \\
	( pendingCount = pendingOne \land pendingCount' = pendingTwo \land
	pendingDeltaOne' = pendingDeltaOne \land pendingDeltaTwo' = delta? \land
	pendingDeltaThree' = pendingDeltaThree \land pendingDeltaFour' = pendingDeltaFour \land
	pendingDeltaFive' = pendingDeltaFive \land pendingDeltaSix' = pendingDeltaSix ) \lor \\
	( pendingCount = pendingTwo \land pendingCount' = pendingThree \land
	pendingDeltaOne' = pendingDeltaOne \land pendingDeltaTwo' = pendingDeltaTwo \land
	pendingDeltaThree' = delta? \land pendingDeltaFour' = pendingDeltaFour \land
	pendingDeltaFive' = pendingDeltaFive \land pendingDeltaSix' = pendingDeltaSix ) \lor \\
	( pendingCount = pendingThree \land pendingCount' = pendingFour \land
	pendingDeltaOne' = pendingDeltaOne \land pendingDeltaTwo' = pendingDeltaTwo \land
	pendingDeltaThree' = pendingDeltaThree \land pendingDeltaFour' = delta? \land
	pendingDeltaFive' = pendingDeltaFive \land pendingDeltaSix' = pendingDeltaSix ) \lor \\
	( pendingCount = pendingFour \land pendingCount' = pendingFive \land
	pendingDeltaOne' = pendingDeltaOne \land pendingDeltaTwo' = pendingDeltaTwo \land
	pendingDeltaThree' = pendingDeltaThree \land pendingDeltaFour' = pendingDeltaFour \land
	pendingDeltaFive' = delta? \land pendingDeltaSix' = pendingDeltaSix ) \lor \\
	( pendingCount = pendingFive \land pendingCount' = pendingSix \land
	pendingDeltaOne' = pendingDeltaOne \land pendingDeltaTwo' = pendingDeltaTwo \land
	pendingDeltaThree' = pendingDeltaThree \land pendingDeltaFour' = pendingDeltaFour \land
	pendingDeltaFive' = pendingDeltaFive \land pendingDeltaSix' = delta? )
\end{schema}
% POINTS-CHECK END effect.append

% POINTS-CHECK BEGIN effect.checkpoint definition
\begin{schema}{DirectoryCheckpoint}
	\Delta TxnVars
	\where
	visibleNS' = visibleNS \land durableNS' = visibleNS \land
	pendingCount' = pendingZero \\
	pendingDeltaOne' = pendingDeltaOne \land pendingDeltaTwo' = pendingDeltaTwo \land
	pendingDeltaThree' = pendingDeltaThree \land
	pendingDeltaFour' = pendingDeltaFour \land
	pendingDeltaFive' = pendingDeltaFive \land pendingDeltaSix' = pendingDeltaSix \\
	namespaceMutated' = namespaceMutated \\
	cleanupDue' = \{ p : \dom cleanupDue |
	\lnot ( \{ p \mapsto file~(cleanupDue~p) \} \subseteq visibleNS ) @ p \}
	\ndres cleanupDue \\
	cleanupPending' = \{ p : \dom cleanupPending |
	\lnot ( \{ p \mapsto file~(cleanupDue~p) \} \subseteq visibleNS ) @ p \}
	\ndres cleanupPending \\
	owned' = owned \land retainedResults' = retainedResults \land
	cleanupNotes' = cleanupNotes \\
	( committed = yes \land
	\{ configPath \mapsto file~candidateID \} \subseteq visibleNS
	\implies configSyncAck' = yes ) \\
	( committed = no \lor
	\lnot ( \{ configPath \mapsto file~candidateID \} \subseteq visibleNS )
	\implies configSyncAck' = configSyncAck ) \\
	( backupAbsenceConfirmed' = yes \iff backupAbsenceConfirmed = yes \lor \\
	\t1 ( ownedBackupPublication \neq \emptyset \land \\
	\t2 ( \forall p : \dom ownedBackupPublication @ \\
	\t3 \lnot ( \{ p \mapsto file~(ownedBackupPublication~p) \}
	\subseteq visibleNS ) ) ) )
\end{schema}
% POINTS-CHECK END effect.checkpoint

Every operation below declares $\Delta TxnVars$, uses the exact frame schemas
for untouched groups, and assigns every member of every changed group.  This is
the authoritative relation; extraction never invents a frame.  Operations that
append a pending delta or establish a directory checkpoint declaration-include
the corresponding authoritative effect schema, rather than mentioning it only
as a predicate.

% POINTS-CHECK BEGIN operation.candidatefailure operation
\begin{schema}{CandidateFailure}
	\Delta TxnVars \\
	TxnStaticFrame \\
	TxnNamespaceFrame \\
	TxnImageFrame \\
	TxnStageFrame \\
	TxnPublicationFrame \\
	TxnOwnershipFrame \\
	TxnAttemptFrame \\
	TxnBackupHistoryFrame
	\where
	life = running \land reportState = reportPending \land lockHeld = yes \land
	cancelled = no \land phase = staging \land stageIDs = \emptyset \land owned = \emptyset \\
	phase' = aborting \land committed' = committed \land
	configSyncAck' = configSyncAck \land cancelled' = cancelled \land
	lockHeld' = lockHeld \land backupProgress' = backupProgress \land
	life' = life \land reportState' = reportState \\
	events' = events \cup \{ wCandidate \}
\end{schema}
% POINTS-CHECK END operation.candidatefailure

% POINTS-CHECK BEGIN operation.createstage operation
\begin{schema}{StageRole}
	TxnVars \\
	p? : Path; id? : Identity
	\where
	( p? = configStage \land id? = candidateID ) \lor \\
	( p? = backupStage \land id? \in backupIDs \land conflicts \neq \emptyset \land
	candidateID \in closeOK )
\end{schema}

\begin{schema}{CreateReady}
	TxnVars \\
	p? : Path; id? : Identity
	\where
	life = running \land reportState = reportPending \land lockHeld = yes \land
	cancelled = no \land phase = staging \\
	StageRole \\
	p? \notin \dom visibleNS \cup \dom owned \\
	id? \in \dom intendedImages \land id? \notin refs~visibleNS \cup refs~durableNS
	\cup \ran owned \cup \dom visibleImages \cup \dom syncedImages
	\cup \dom dirtyChoices \cup stageIDs \cup \dom publicationImages \\
	( id? \in backupIDs \implies id? \in \dom snapshotImages \land
	intendedImages~id? = snapshotImages~id? )
\end{schema}

\begin{schema}{CreateStage}
	\Delta TxnVars \\
	TxnStaticFrame \\
	TxnPublicationFrame \\
	TxnControlFrame \\
	TxnAttemptFrame \\
	TxnEventFrame \\
	TxnBackupHistoryFrame \\
	AppendPendingDelta \\
	p? : Path; id? : Identity; initialImage? : FileImage; delta? : Delta
	\where
	CreateReady \\
	imageBytes~initialImage? = emptyBytes \land
	imageMode~initialImage? = imageMode~(intendedImages~id?) \land
	imageOwner~initialImage? = imageOwner~(intendedImages~id?) \\
	( imageState~initialImage? = complete \iff
	imageBytes~initialImage? = imageBytes~(intendedImages~id?) ) \\
	deltaRemove~delta? = \emptyset \land
	deltaPut~delta? = \{ p? \mapsto file~id? \} \\
	visibleNS' = (applyDelta~delta?)~visibleNS \land durableNS' = durableNS \land
	namespaceMutated' = yes \\
	visibleImages' = visibleImages \oplus \{ id? \mapsto initialImage? \} \land
	syncedImages' = syncedImages \\
	dirtyChoices' = dirtyChoices \oplus \{ id? \mapsto \{ i : FileImage | \\
	\t1 imageBytes~i \mapsto imageBytes~initialImage? \in bytePrefix \land
	imageMode~i = imageMode~(intendedImages~id?) \land \\
	\t1 imageOwner~i = imageOwner~(intendedImages~id?) \land \\
	\t1 ( imageState~i = complete \iff
	imageBytes~i = imageBytes~(intendedImages~id?) ) @ i \} \} \\
	stageIDs' = stageIDs \cup \{ id? \} \land sealedIDs' = sealedIDs \land
	closeOK' = closeOK \\
	owned' = owned \oplus \{ p? \mapsto id? \} \land
	cleanupDue' = cleanupDue \oplus \{ p? \mapsto id? \} \land
	cleanupPending' = cleanupPending \land retainedResults' = retainedResults \land
	cleanupNotes' = cleanupNotes
\end{schema}
% POINTS-CHECK END operation.createstage

% POINTS-CHECK BEGIN operation.createstagefailure operation
\begin{schema}{CreateStageFailure}
	\Delta TxnVars \\
	TxnStaticFrame \\
	TxnNamespaceFrame \\
	TxnImageFrame \\
	TxnStageFrame \\
	TxnPublicationFrame \\
	TxnOwnershipFrame \\
	TxnAttemptFrame \\
	TxnBackupHistoryFrame \\
	p? : Path; id? : Identity
	\where
	CreateReady \\
	phase' = aborting \land committed' = committed \land
	configSyncAck' = configSyncAck \land cancelled' = cancelled \land
	lockHeld' = lockHeld \land backupProgress' = backupProgress \land
	life' = life \land reportState' = reportState \\
	( id? = candidateID \implies events' = events \cup \{ wConfigStage \} ) \\
	( id? \in backupIDs \implies events' = events \cup \{ wBackupStage \} )
\end{schema}
% POINTS-CHECK END operation.createstagefailure

% POINTS-CHECK BEGIN operation.writestage operation
\begin{schema}{WriteStage}
	\Delta TxnVars \\
	TxnStaticFrame \\
	TxnNamespaceFrame \\
	TxnStageFrame \\
	TxnPublicationFrame \\
	TxnOwnershipFrame \\
	TxnControlFrame \\
	TxnAttemptFrame \\
	TxnEventFrame \\
	TxnBackupHistoryFrame \\
	id? : Identity; image? : FileImage
	\where
	life = running \land reportState = reportPending \land lockHeld = yes \land
	cancelled = no \land phase = staging \\
	id? \in stageIDs \cap \dom visibleImages \cap \dom intendedImages \land
	id? \notin sealedIDs \cup closeOK \\
	imageBytes~(visibleImages~id?) \mapsto imageBytes~image? \in bytePrefix \land
	imageBytes~(visibleImages~id?) \neq imageBytes~image? \\
	imageBytes~image? \mapsto imageBytes~(intendedImages~id?) \in bytePrefix \\
	imageMode~image? = imageMode~(intendedImages~id?) \land
	imageOwner~image? = imageOwner~(intendedImages~id?) \\
	( imageState~image? = complete \iff
	imageBytes~image? = imageBytes~(intendedImages~id?) ) \\
	visibleImages' = visibleImages \oplus \{ id? \mapsto image? \} \land
	syncedImages' = syncedImages \\
	dirtyChoices' = dirtyChoices \oplus \{ id? \mapsto \{ i : FileImage | \\
	\t1 imageBytes~i \mapsto imageBytes~image? \in bytePrefix \land
	imageMode~i = imageMode~(intendedImages~id?) \land \\
	\t1 imageOwner~i = imageOwner~(intendedImages~id?) \land \\
	\t1 ( imageState~i = complete \iff
	imageBytes~i = imageBytes~(intendedImages~id?) ) @ i \} \}
\end{schema}
% POINTS-CHECK END operation.writestage

% POINTS-CHECK BEGIN operation.writestagefailure operation
\begin{schema}{WriteStageFailure}
	\Delta TxnVars \\
	TxnStaticFrame \\
	TxnNamespaceFrame \\
	TxnImageFrame \\
	TxnStageFrame \\
	TxnPublicationFrame \\
	TxnOwnershipFrame \\
	TxnAttemptFrame \\
	TxnBackupHistoryFrame \\
	id? : Identity
	\where
	life = running \land reportState = reportPending \land lockHeld = yes \land
	cancelled = no \land phase = staging \land
	id? \in stageIDs \setminus ( sealedIDs \cup closeOK ) \\
	phase' = aborting \land committed' = committed \land
	configSyncAck' = configSyncAck \land cancelled' = cancelled \land
	lockHeld' = lockHeld \land backupProgress' = backupProgress \land
	life' = life \land reportState' = reportState \\
	( id? = candidateID \implies events' = events \cup \{ wConfigStage \} ) \\
	( id? \in backupIDs \implies events' = events \cup \{ wBackupStage \} )
\end{schema}
% POINTS-CHECK END operation.writestagefailure

% POINTS-CHECK BEGIN operation.filesyncstage operation
\begin{schema}{FileSyncStage}
	\Delta TxnVars \\
	TxnStaticFrame \\
	TxnNamespaceFrame \\
	TxnPublicationFrame \\
	TxnOwnershipFrame \\
	TxnControlFrame \\
	TxnAttemptFrame \\
	TxnEventFrame \\
	TxnBackupHistoryFrame \\
	id? : Identity
	\where
	life = running \land reportState = reportPending \land lockHeld = yes \land
	cancelled = no \land phase = staging \\
	WrittenComplete \land id? \in stageIDs \setminus ( sealedIDs \cup closeOK ) \\
	visibleImages' = visibleImages \land
	syncedImages' = syncedImages \oplus \{ id? \mapsto intendedImages~id? \} \land
	dirtyChoices' = \{ id? \} \ndres dirtyChoices \\
	stageIDs' = stageIDs \land sealedIDs' = sealedIDs \cup \{ id? \} \land
	closeOK' = closeOK
\end{schema}
% POINTS-CHECK END operation.filesyncstage

% POINTS-CHECK BEGIN operation.filesyncstagefailure operation
\begin{schema}{FileSyncStageFailure}
	WriteStageFailure
	\where WrittenComplete
\end{schema}
% POINTS-CHECK END operation.filesyncstagefailure

% POINTS-CHECK BEGIN operation.closestage operation
\begin{schema}{CloseStage}
	\Delta TxnVars \\
	TxnStaticFrame \\
	TxnNamespaceFrame \\
	TxnImageFrame \\
	TxnPublicationFrame \\
	TxnOwnershipFrame \\
	TxnAttemptFrame \\
	TxnEventFrame \\
	TxnBackupHistoryFrame \\
	id? : Identity
	\where
	life = running \land reportState = reportPending \land lockHeld = yes \land
	cancelled = no \land phase = staging \\
	id? \in stageIDs \setminus closeOK \\
	WrittenComplete \\
	% POINTS-MUTATION BEGIN m2.close.file-sync
	FileSynced \\
	% POINTS-MUTATION END m2.close.file-sync
	stageIDs' = stageIDs \land sealedIDs' = sealedIDs \land
	closeOK' = closeOK \cup \{ id? \} \\
	phase' = phase \land committed' = committed \land
	configSyncAck' = configSyncAck \land cancelled' = cancelled \land
	lockHeld' = lockHeld \land life' = life \land reportState' = reportState \\
	( id? \in backupIDs \implies backupProgress' = backupStaged ) \\
	( id? \notin backupIDs \implies backupProgress' = backupProgress )
\end{schema}
% POINTS-CHECK END operation.closestage

% POINTS-CHECK BEGIN operation.closestagefailure operation
\begin{schema}{CloseStageFailure}
	\Delta TxnVars \\
	TxnStaticFrame \\
	TxnNamespaceFrame \\
	TxnImageFrame \\
	TxnStageFrame \\
	TxnPublicationFrame \\
	TxnOwnershipFrame \\
	TxnAttemptFrame \\
	TxnBackupHistoryFrame \\
	id? : Identity
	\where
	life = running \land reportState = reportPending \land lockHeld = yes \land
	cancelled = no \land phase = staging \\
	id? \in stageIDs \setminus closeOK \\
	WrittenComplete \\
	FileSynced \\
	phase' = aborting \land committed' = committed \land
	configSyncAck' = configSyncAck \land cancelled' = cancelled \land
	lockHeld' = lockHeld \land backupProgress' = backupProgress \land
	life' = life \land reportState' = reportState \\
	( id? = candidateID \implies events' = events \cup \{ wConfigStage \} ) \\
	( id? \in backupIDs \implies events' = events \cup \{ wBackupStage \} )
\end{schema}
% POINTS-CHECK END operation.closestagefailure

% POINTS-CHECK BEGIN operation.beginrevalidation operation
\begin{schema}{BeginRevalidation}
	\Delta TxnVars \\
	TxnStaticFrame \\
	TxnNamespaceFrame \\
	TxnImageFrame \\
	TxnStageFrame \\
	TxnPublicationFrame \\
	TxnOwnershipFrame \\
	TxnAttemptFrame \\
	TxnEventFrame \\
	TxnBackupHistoryFrame
	\where
	life = running \land reportState = reportPending \land lockHeld = yes \land
	cancelled = no \land phase = staging \\
	( \exists id? : Identity @
	id? = candidateID \land WrittenComplete ) \land
	candidateID \in closeOK \\
	BackupBarrier \\
	phase' = revalidating \land committed' = committed \land
	configSyncAck' = configSyncAck \land cancelled' = cancelled \land
	lockHeld' = lockHeld \land backupProgress' = backupProgress \land
	life' = life \land reportState' = reportState
\end{schema}
% POINTS-CHECK END operation.beginrevalidation

% POINTS-CHECK BEGIN operation.revalidationsuccess operation
\begin{schema}{RevalidationSuccess}
	\Delta TxnVars \\
	TxnStaticFrame \\
	TxnNamespaceFrame \\
	TxnImageFrame \\
	TxnStageFrame \\
	TxnPublicationFrame \\
	TxnOwnershipFrame \\
	TxnEventFrame \\
	TxnBackupHistoryFrame
	\where
	life = running \land reportState = reportPending \land lockHeld = yes \land
	cancelled = no \land phase = revalidating \land SnapshotMatches \\
	phase' = publishingConfig \land committed' = committed \land
	configSyncAck' = configSyncAck \land cancelled' = cancelled \land
	lockHeld' = lockHeld \land backupProgress' = backupProgress \land
	life' = life \land reportState' = reportState \\
	revalidationPassed' = yes \land confirmationAttempted' = confirmationAttempted \land
	cleanupWork' = cleanupWork \land cleanupTried' = cleanupTried \land
	cleanupSyncAttempted' = cleanupSyncAttempted \land
	releaseAttempted' = releaseAttempted
\end{schema}
% POINTS-CHECK END operation.revalidationsuccess

% POINTS-CHECK BEGIN operation.revalidationiofailure operation
\begin{schema}{RevalidationIOFailure}
	\Delta TxnVars \\
	TxnStaticFrame \\
	TxnNamespaceFrame \\
	TxnImageFrame \\
	TxnStageFrame \\
	TxnPublicationFrame \\
	TxnOwnershipFrame \\
	TxnAttemptFrame \\
	TxnBackupHistoryFrame
	\where
	life = running \land reportState = reportPending \land lockHeld = yes \land
	cancelled = no \land phase = revalidating \\
	phase' = aborting \land committed' = committed \land
	configSyncAck' = configSyncAck \land cancelled' = cancelled \land
	lockHeld' = lockHeld \land backupProgress' = backupProgress \land
	life' = life \land reportState' = reportState \\
	events' = events \cup \{ wChanged \}
\end{schema}
% POINTS-CHECK END operation.revalidationiofailure

% POINTS-CHECK BEGIN operation.commitabsent operation
\begin{schema}{CommitAbsent}
	\Delta TxnVars \\
	TxnStaticFrame \\
	TxnImageFrame \\
	TxnStageFrame \\
	TxnAttemptFrame \\
	TxnEventFrame \\
	TxnBackupHistoryFrame \\
	AppendPendingDelta \\
	delta? : Delta
	\where
	life = running \land reportState = reportPending \land lockHeld = yes \land
	cancelled = no \\
	% POINTS-MUTATION BEGIN m1.commit-absent.revalidation
	RevalidationBarrier \\
	% POINTS-MUTATION END m1.commit-absent.revalidation
	SnapshotMatches \land BackupBarrier \\
	expectedConfig = \emptyset \land configPath \notin \dom visibleNS \\
	configStage \in \dom cleanupDue \land cleanupDue~configStage = candidateID \land
	\{ configStage \mapsto file~candidateID \} \subseteq visibleNS \\
	( \exists id? : Identity @ id? = candidateID \land PublicationQualified ) \\
	deltaRemove~delta? = \emptyset \land
	deltaPut~delta? = \{ configPath \mapsto file~candidateID \} \\
	visibleNS' = (applyDelta~delta?)~visibleNS \land durableNS' = durableNS \land
	namespaceMutated' = yes \\
	published' = published \cup \{ configPath \mapsto candidateID \} \land
	publicationImages' = publicationImages \oplus
	\{ candidateID \mapsto intendedImages~candidateID \} \\
	owned' = owned \oplus \{ configPath \mapsto candidateID \} \land
	cleanupDue' = \{ backupPath \} \ndres cleanupDue \land
	cleanupPending' = cleanupPending \land
	retainedResults' = retainedResults \oplus \{ configPath \mapsto candidateID \}
	\oplus ( \{ backupPath \} \dres owned ) \land cleanupNotes' = cleanupNotes \\
	phase' = confirming \land committed' = yes \land configSyncAck' = no \land
	cancelled' = cancelled \land lockHeld' = lockHeld \land
	backupProgress' = backupProgress \land life' = life \land
	reportState' = reportState
\end{schema}
% POINTS-CHECK END operation.commitabsent

% POINTS-CHECK BEGIN operation.commitabsentfailure operation
\begin{schema}{CommitAbsentFailure}
	\Delta TxnVars \\
	TxnStaticFrame \\
	TxnNamespaceFrame \\
	TxnImageFrame \\
	TxnStageFrame \\
	TxnPublicationFrame \\
	TxnOwnershipFrame \\
	TxnAttemptFrame \\
	TxnBackupHistoryFrame
	\where
	life = running \land reportState = reportPending \land lockHeld = yes \land
	cancelled = no \\
	RevalidationBarrier \\
	SnapshotMatches \land BackupBarrier \\
	expectedConfig = \emptyset \land configPath \notin \dom visibleNS \\
	configStage \in \dom cleanupDue \land cleanupDue~configStage = candidateID \land
	\{ configStage \mapsto file~candidateID \} \subseteq visibleNS \\
	( \exists id? : Identity @ id? = candidateID \land PublicationQualified ) \\
	phase' = aborting \land committed' = committed \land
	configSyncAck' = configSyncAck \land cancelled' = cancelled \land
	lockHeld' = lockHeld \land backupProgress' = backupProgress \land
	life' = life \land reportState' = reportState \\
	events' = events \cup \{ wConfigPublish \}
\end{schema}
% POINTS-CHECK END operation.commitabsentfailure

% POINTS-CHECK BEGIN operation.confirmconfigsuccess operation
\begin{schema}{ConfirmConfigSuccess}
	\Delta TxnVars \\
	TxnStaticFrame \\
	TxnImageFrame \\
	TxnStageFrame \\
	TxnPublicationFrame \\
	TxnEventFrame \\
	DirectoryCheckpoint
	\where
	life = running \land reportState = reportPending \land lockHeld = yes \land
	phase = confirming \land committed = yes \land confirmationAttempted = no \\
	phase' = cleaning \land committed' = committed \land cancelled' = cancelled \land
	lockHeld' = lockHeld \land backupProgress' = backupProgress \land
	life' = life \land reportState' = reportState \\
	revalidationPassed' = revalidationPassed \land confirmationAttempted' = yes \land
	cleanupWork' = cleanupWork \land cleanupTried' = cleanupTried \land
	cleanupSyncAttempted' = cleanupSyncAttempted \land releaseAttempted' = releaseAttempted \\
	ownedBackupPublication' = ownedBackupPublication
\end{schema}
% POINTS-CHECK END operation.confirmconfigsuccess

% POINTS-CHECK BEGIN operation.confirmconfigfailure operation
\begin{schema}{ConfirmConfigFailure}
	\Delta TxnVars \\
	TxnStaticFrame \\
	TxnNamespaceFrame \\
	TxnImageFrame \\
	TxnStageFrame \\
	TxnPublicationFrame \\
	TxnOwnershipFrame \\
	TxnBackupHistoryFrame
	\where
	life = running \land reportState = reportPending \land lockHeld = yes \land
	phase = confirming \land committed = yes \land confirmationAttempted = no \\
	phase' = cleaning \land committed' = committed \land
	configSyncAck' = configSyncAck \land cancelled' = cancelled \land
	lockHeld' = lockHeld \land backupProgress' = backupProgress \land
	life' = life \land reportState' = reportState \\
	revalidationPassed' = revalidationPassed \land confirmationAttempted' = yes \land
	cleanupWork' = cleanupWork \land cleanupTried' = cleanupTried \land
	cleanupSyncAttempted' = cleanupSyncAttempted \land releaseAttempted' = releaseAttempted \\
	events' = events \cup \{ dConfigSync \}
\end{schema}
% POINTS-CHECK END operation.confirmconfigfailure

% POINTS-CHECK BEGIN operation.begincleanup operation
\begin{schema}{BeginCleanup}
	\Delta TxnVars \\
	TxnStaticFrame \\
	TxnNamespaceFrame \\
	TxnImageFrame \\
	TxnStageFrame \\
	TxnPublicationFrame \\
	TxnOwnershipFrame \\
	TxnEventFrame \\
	TxnBackupHistoryFrame
	\where
	life = running \land reportState = reportPending \land lockHeld = yes \land
	phase \in \{ aborting, cleaning \} \\
	phase' = sweeping \land committed' = committed \land
	configSyncAck' = configSyncAck \land cancelled' = cancelled \land
	lockHeld' = lockHeld \land backupProgress' = backupProgress \land
	life' = life \land reportState' = reportState \\
	revalidationPassed' = revalidationPassed \land
	confirmationAttempted' = confirmationAttempted \land
	cleanupWork' = cleanupDue \land cleanupTried' = \emptyset \land
	cleanupSyncAttempted' = cleanupSyncAttempted \land
	releaseAttempted' = releaseAttempted
\end{schema}
% POINTS-CHECK END operation.begincleanup

% POINTS-CHECK BEGIN operation.unlinkowned operation
\begin{schema}{UnlinkOwned}
	\Delta TxnVars \\
	TxnStaticFrame \\
	TxnImageFrame \\
	TxnStageFrame \\
	TxnPublicationFrame \\
	TxnEventFrame \\
	TxnBackupHistoryFrame \\
	AppendPendingDelta \\
	p? : Path; delta? : Delta
	\where
	life = running \land reportState = reportPending \land lockHeld = yes \land
	phase = sweeping \land p? \in \dom cleanupWork \setminus cleanupTried \land
	p? \notin \dom retainedResults \\
	\{ p? \mapsto file~(cleanupWork~p?) \} \subseteq visibleNS \\
	deltaRemove~delta? = \{ p? \} \land deltaPut~delta? = \emptyset \\
	visibleNS' = (applyDelta~delta?)~visibleNS \land durableNS' = durableNS \land
	namespaceMutated' = yes \\
	owned' = owned \land cleanupDue' = cleanupDue \land
	cleanupPending' = cleanupPending \oplus \{ p? \mapsto cleanupWork~p? \} \land
	retainedResults' = retainedResults \land cleanupNotes' = cleanupNotes \\
	phase' = phase \land committed' = committed \land
	configSyncAck' = configSyncAck \land cancelled' = cancelled \land
	lockHeld' = lockHeld \land backupProgress' = backupProgress \land
	life' = life \land reportState' = reportState \\
	revalidationPassed' = revalidationPassed \land
	confirmationAttempted' = confirmationAttempted \land
	cleanupWork' = cleanupWork \land cleanupTried' = cleanupTried \cup \{ p? \} \land
	cleanupSyncAttempted' = cleanupSyncAttempted \land releaseAttempted' = releaseAttempted
\end{schema}
% POINTS-CHECK END operation.unlinkowned

% POINTS-CHECK BEGIN operation.cleanupalreadyabsent operation
\begin{schema}{CleanupAlreadyAbsent}
	\Delta TxnVars \\
	TxnStaticFrame \\
	TxnNamespaceFrame \\
	TxnImageFrame \\
	TxnStageFrame \\
	TxnPublicationFrame \\
	TxnControlFrame \\
	TxnEventFrame \\
	TxnBackupHistoryFrame \\
	p? : Path
	\where
	life = running \land reportState = reportPending \land lockHeld = yes \land
	phase = sweeping \land p? \in \dom cleanupWork \setminus cleanupTried \land
	p? \notin \dom retainedResults \cup \dom visibleNS \\
	owned' = owned \land cleanupDue' = cleanupDue \land
	cleanupPending' = cleanupPending \oplus \{ p? \mapsto cleanupWork~p? \} \land
	retainedResults' = retainedResults \land cleanupNotes' = cleanupNotes \\
	revalidationPassed' = revalidationPassed \land
	confirmationAttempted' = confirmationAttempted \land
	cleanupWork' = cleanupWork \land cleanupTried' = cleanupTried \cup \{ p? \} \land
	cleanupSyncAttempted' = cleanupSyncAttempted \land releaseAttempted' = releaseAttempted
\end{schema}
% POINTS-CHECK END operation.cleanupalreadyabsent

% POINTS-CHECK BEGIN operation.cleanupunlinkfailed operation
\begin{schema}{CleanupUnlinkFailed}
	\Delta TxnVars \\
	TxnStaticFrame \\
	TxnNamespaceFrame \\
	TxnImageFrame \\
	TxnStageFrame \\
	TxnPublicationFrame \\
	TxnControlFrame \\
	TxnBackupHistoryFrame \\
	p? : Path
	\where
	life = running \land reportState = reportPending \land lockHeld = yes \land
	phase = sweeping \land p? \in \dom cleanupWork \setminus cleanupTried \land
	p? \notin \dom retainedResults \land
	\{ p? \mapsto file~(cleanupWork~p?) \} \subseteq visibleNS \\
	owned' = owned \land cleanupDue' = cleanupDue \land
	cleanupPending' = cleanupPending \land retainedResults' = retainedResults \land
	cleanupNotes' = cleanupNotes \oplus \{ p? \mapsto unlinkFailed \} \\
	revalidationPassed' = revalidationPassed \land
	confirmationAttempted' = confirmationAttempted \land
	cleanupWork' = cleanupWork \land cleanupTried' = cleanupTried \cup \{ p? \} \land
	cleanupSyncAttempted' = cleanupSyncAttempted \land releaseAttempted' = releaseAttempted \\
	events' = events \cup \{ cUnlink \}
\end{schema}
% POINTS-CHECK END operation.cleanupunlinkfailed

% POINTS-CHECK BEGIN operation.cleanupinspectionfailed operation
\begin{schema}{CleanupInspectionFailed}
	\Delta TxnVars \\
	TxnStaticFrame \\
	TxnNamespaceFrame \\
	TxnImageFrame \\
	TxnStageFrame \\
	TxnPublicationFrame \\
	TxnControlFrame \\
	TxnBackupHistoryFrame \\
	p? : Path
	\where
	life = running \land reportState = reportPending \land lockHeld = yes \land
	phase = sweeping \land p? \in \dom cleanupWork \setminus cleanupTried \land
	p? \notin \dom retainedResults \\
	owned' = owned \land cleanupDue' = cleanupDue \land
	cleanupPending' = cleanupPending \land retainedResults' = retainedResults \land
	cleanupNotes' = cleanupNotes \oplus \{ p? \mapsto inspectFailed \} \\
	revalidationPassed' = revalidationPassed \land
	confirmationAttempted' = confirmationAttempted \land
	cleanupWork' = cleanupWork \land cleanupTried' = cleanupTried \cup \{ p? \} \land
	cleanupSyncAttempted' = cleanupSyncAttempted \land releaseAttempted' = releaseAttempted \\
	events' = events \cup \{ cInspect \}
\end{schema}
% POINTS-CHECK END operation.cleanupinspectionfailed

% POINTS-CHECK BEGIN operation.endcleanupsweep operation
\begin{schema}{EndCleanupSweep}
	\Delta TxnVars \\
	TxnStaticFrame \\
	TxnNamespaceFrame \\
	TxnImageFrame \\
	TxnStageFrame \\
	TxnPublicationFrame \\
	TxnOwnershipFrame \\
	TxnAttemptFrame \\
	TxnEventFrame \\
	TxnBackupHistoryFrame
	\where
	life = running \land reportState = reportPending \land lockHeld = yes \land
	phase = sweeping \land cleanupTried = \dom cleanupWork \\
	phase' = cleanupSync \land committed' = committed \land
	configSyncAck' = configSyncAck \land cancelled' = cancelled \land
	lockHeld' = lockHeld \land backupProgress' = backupProgress \land
	life' = life \land reportState' = reportState
\end{schema}
% POINTS-CHECK END operation.endcleanupsweep

% POINTS-CHECK BEGIN operation.cleanupsyncsuccess operation
\begin{schema}{CleanupSyncSuccess}
	\Delta TxnVars \\
	TxnStaticFrame \\
	TxnImageFrame \\
	TxnStageFrame \\
	TxnPublicationFrame \\
	TxnEventFrame \\
	DirectoryCheckpoint
	\where
	life = running \land reportState = reportPending \land lockHeld = yes \land
	phase = cleanupSync \land cleanupSyncAttempted = no \\
	phase' = releasing \land committed' = committed \land cancelled' = cancelled \land
	lockHeld' = lockHeld \land backupProgress' = backupProgress \land
	life' = life \land reportState' = reportState \\
	revalidationPassed' = revalidationPassed \land
	confirmationAttempted' = confirmationAttempted \land
	cleanupWork' = cleanupWork \land cleanupTried' = cleanupTried \land
	cleanupSyncAttempted' = yes \land releaseAttempted' = releaseAttempted \\
	ownedBackupPublication' = ownedBackupPublication
\end{schema}
% POINTS-CHECK END operation.cleanupsyncsuccess

% POINTS-CHECK BEGIN operation.cleanupsyncfailure operation
\begin{schema}{CleanupSyncFailure}
	\Delta TxnVars \\
	TxnStaticFrame \\
	TxnNamespaceFrame \\
	TxnImageFrame \\
	TxnStageFrame \\
	TxnPublicationFrame \\
	TxnOwnershipFrame \\
	TxnBackupHistoryFrame
	\where
	life = running \land reportState = reportPending \land lockHeld = yes \land
	phase = cleanupSync \land cleanupSyncAttempted = no \\
	phase' = releasing \land committed' = committed \land
	configSyncAck' = configSyncAck \land cancelled' = cancelled \land
	lockHeld' = lockHeld \land backupProgress' = backupProgress \land
	life' = life \land reportState' = reportState \\
	revalidationPassed' = revalidationPassed \land
	confirmationAttempted' = confirmationAttempted \land
	cleanupWork' = cleanupWork \land cleanupTried' = cleanupTried \land
	cleanupSyncAttempted' = yes \land releaseAttempted' = releaseAttempted \\
	events' = events \cup \{ cDeleteSync \}
\end{schema}
% POINTS-CHECK END operation.cleanupsyncfailure

% POINTS-CHECK BEGIN operation.releasesuccess operation
\begin{schema}{ReleaseSuccess}
	\Delta TxnVars \\
	TxnStaticFrame \\
	TxnNamespaceFrame \\
	TxnImageFrame \\
	TxnStageFrame \\
	TxnPublicationFrame \\
	TxnOwnershipFrame \\
	TxnEventFrame \\
	TxnBackupHistoryFrame
	\where
	life = running \land reportState = reportPending \land phase = releasing \land
	releaseAttempted = no \\
	phase' = reporting \land committed' = committed \land
	configSyncAck' = configSyncAck \land cancelled' = cancelled \land
	lockHeld' = no \land backupProgress' = backupProgress \land
	life' = life \land reportState' = reportState \\
	revalidationPassed' = revalidationPassed \land
	confirmationAttempted' = confirmationAttempted \land
	cleanupWork' = cleanupWork \land cleanupTried' = cleanupTried \land
	cleanupSyncAttempted' = cleanupSyncAttempted \land releaseAttempted' = yes
\end{schema}
% POINTS-CHECK END operation.releasesuccess

% POINTS-CHECK BEGIN operation.releasefailure operation
\begin{schema}{ReleaseFailure}
	\Delta TxnVars \\
	TxnStaticFrame \\
	TxnNamespaceFrame \\
	TxnImageFrame \\
	TxnStageFrame \\
	TxnPublicationFrame \\
	TxnOwnershipFrame \\
	TxnBackupHistoryFrame
	\where
	life = running \land reportState = reportPending \land phase = releasing \land
	releaseAttempted = no \\
	phase' = reporting \land committed' = committed \land
	configSyncAck' = configSyncAck \land cancelled' = cancelled \land
	lockHeld' = no \land backupProgress' = backupProgress \land
	life' = life \land reportState' = reportState \\
	revalidationPassed' = revalidationPassed \land
	confirmationAttempted' = confirmationAttempted \land
	cleanupWork' = cleanupWork \land cleanupTried' = cleanupTried \land
	cleanupSyncAttempted' = cleanupSyncAttempted \land releaseAttempted' = yes \\
	events' = events \cup \{ cRelease \}
\end{schema}
% POINTS-CHECK END operation.releasefailure

% POINTS-CHECK BEGIN operation.reportdelivered operation
\begin{schema}{ReportDelivered}
	\Delta TxnVars \\
	TxnStaticFrame \\
	TxnNamespaceFrame \\
	TxnImageFrame \\
	TxnStageFrame \\
	TxnPublicationFrame \\
	TxnOwnershipFrame \\
	TxnAttemptFrame \\
	TxnEventFrame \\
	TxnBackupHistoryFrame
	\where
	life = running \land phase = reporting \land reportState = reportPending \land
	releaseAttempted = yes \land lockHeld = no \\
	phase' = phase \land committed' = committed \land
	configSyncAck' = configSyncAck \land cancelled' = cancelled \land
	lockHeld' = lockHeld \land backupProgress' = backupProgress \land
	life' = life \land reportState' = reportDelivered
\end{schema}
% POINTS-CHECK END operation.reportdelivered

% POINTS-CHECK BEGIN operation.reportfailed operation
\begin{schema}{ReportFailed}
	\Delta TxnVars \\
	TxnStaticFrame \\
	TxnNamespaceFrame \\
	TxnImageFrame \\
	TxnStageFrame \\
	TxnPublicationFrame \\
	TxnOwnershipFrame \\
	TxnAttemptFrame \\
	TxnBackupHistoryFrame
	\where
	life = running \land phase = reporting \land reportState = reportPending \land
	releaseAttempted = yes \land lockHeld = no \\
	phase' = phase \land committed' = committed \land
	configSyncAck' = configSyncAck \land cancelled' = cancelled \land
	lockHeld' = lockHeld \land backupProgress' = backupProgress \land
	life' = life \land reportState' = reportFailed \\
	( committed = yes \implies events' = events \cup \{ oReportAfter \} ) \\
	( committed = no \implies events' = events \cup \{ oReportBefore \} )
\end{schema}
% POINTS-CHECK END operation.reportfailed

% POINTS-CHECK BEGIN operation.finishtransaction operation
\begin{schema}{FinishTransaction}
	\Delta TxnVars \\
	TxnStaticFrame \\
	TxnNamespaceFrame \\
	TxnImageFrame \\
	TxnStageFrame \\
	TxnPublicationFrame \\
	TxnOwnershipFrame \\
	TxnEventFrame \\
	TxnBackupHistoryFrame
	\where
	life = running \land phase = reporting \land
	reportState \in \{ reportDelivered, reportFailed \} \land
	releaseAttempted = yes \land cleanupSyncAttempted = yes \land
	cleanupTried = \dom cleanupWork \land
	( committed = yes \implies confirmationAttempted = yes ) \\
	phase' = finished \land committed' = committed \land
	configSyncAck' = configSyncAck \land cancelled' = cancelled \land
	lockHeld' = no \land backupProgress' = backupProgress \land
	life' = ended \land reportState' = reportState \\
	revalidationPassed' = revalidationPassed \land
	confirmationAttempted' = confirmationAttempted \land
	cleanupWork' = cleanupWork \land cleanupTried' = cleanupTried \land
	cleanupSyncAttempted' = cleanupSyncAttempted \land releaseAttempted' = releaseAttempted
\end{schema}
% POINTS-CHECK END operation.finishtransaction

Normal finish removes executable authority only through $life=ended$,
$phase=finished$, and $lockHeld=no$.  It preserves stage evidence, historical
ownership, cleanup obligations and tombstones, retained results, and cleanup
notes so the completed-result projection can still distinguish unconfirmed
deletion after a failed cleanup checkpoint.  Process death and power loss retain
their separate explicit authority-clearing effects.

% POINTS-CHECK BEGIN operation.observecancellation operation
\begin{schema}{ObserveCancellation}
	\Delta TxnVars \\
	TxnStaticFrame \\
	TxnNamespaceFrame \\
	TxnImageFrame \\
	TxnStageFrame \\
	TxnPublicationFrame \\
	TxnOwnershipFrame \\
	TxnAttemptFrame \\
	TxnBackupHistoryFrame
	\where
	life = running \land reportState = reportPending \land cancelled = no \land
	phase \neq finished \\
	committed' = committed \land configSyncAck' = configSyncAck \land
	cancelled' = yes \land lockHeld' = lockHeld \land
	backupProgress' = backupProgress \land life' = life \land
	reportState' = reportState \\
	( committed = yes \implies phase' = phase \land
	events' = events \cup \{ xAfterCommit \} ) \\
	( committed = no \land phase \in
	\{ staging, publishingBackup, revalidating, publishingConfig \} \implies
	phase' = aborting \land events' = events \cup \{ xDuringWrite \} ) \\
	( committed = no \land phase \notin
	\{ staging, publishingBackup, revalidating, publishingConfig \} \implies
	phase' = phase \land events' = events \cup \{ xDuringWrite \} )
\end{schema}
% POINTS-CHECK END operation.observecancellation

% POINTS-CHECK BEGIN operation.processdeath operation
\begin{schema}{ProcessDeath}
	\Delta TxnVars \\
	TxnStaticFrame \\
	TxnNamespaceFrame \\
	TxnImageFrame \\
	TxnPublicationFrame \\
	TxnEventFrame \\
	TxnBackupHistoryFrame
	\where
	life = running \land phase \neq finished \\
	stageIDs' = \emptyset \land sealedIDs' = \emptyset \land closeOK' = \emptyset \\
	owned' = \emptyset \land cleanupDue' = \emptyset \land
	cleanupPending' = \emptyset \land retainedResults' = \emptyset \land
	cleanupNotes' = \emptyset \\
	phase' = finished \land committed' = committed \land
	configSyncAck' = configSyncAck \land cancelled' = cancelled \land
	lockHeld' = no \land backupProgress' = backupProgress \land
	life' = ended \land reportState' = processDead \\
	revalidationPassed' = revalidationPassed \land
	confirmationAttempted' = confirmationAttempted \land
	cleanupWork' = \emptyset \land cleanupTried' = \emptyset \land
	cleanupSyncAttempted' = cleanupSyncAttempted \land releaseAttempted' = releaseAttempted
\end{schema}
% POINTS-CHECK END operation.processdeath

% POINTS-CHECK BEGIN operation.powerloss operation
\begin{schema}{PowerLoss}
	\Delta TxnVars \\
	TxnStaticFrame \\
	TxnPublicationFrame \\
	TxnEventFrame \\
	TxnBackupHistoryFrame \\
	crashNS : Namespace; crashImages : Identity \pfun FileImage \\
	takeOne, takeTwo, takeThree, takeFour, takeFive, takeSix : Bool \\
	afterOne, afterTwo, afterThree, afterFour, afterFive : Namespace
	\where
	life = running \land phase \neq finished \\
	( pendingCount = pendingZero \implies takeOne = no ) \\
	( pendingCount \notin \{ pendingTwo, pendingThree, pendingFour, pendingFive,
	pendingSix \} \implies takeTwo = no ) \\
	( pendingCount \notin \{ pendingThree, pendingFour, pendingFive, pendingSix \}
	\implies takeThree = no ) \\
	( pendingCount \notin \{ pendingFour, pendingFive, pendingSix \}
	\implies takeFour = no ) \\
	( pendingCount \notin \{ pendingFive, pendingSix \} \implies takeFive = no ) \\
	( pendingCount \neq pendingSix \implies takeSix = no ) \\
	( takeOne = yes \implies afterOne = (applyDelta~pendingDeltaOne)~durableNS ) \land
	( takeOne = no \implies afterOne = durableNS ) \\
	( takeTwo = yes \implies afterTwo = (applyDelta~pendingDeltaTwo)~afterOne ) \land
	( takeTwo = no \implies afterTwo = afterOne ) \\
	( takeThree = yes \implies afterThree = (applyDelta~pendingDeltaThree)~afterTwo )
	\land ( takeThree = no \implies afterThree = afterTwo ) \\
	( takeFour = yes \implies afterFour = (applyDelta~pendingDeltaFour)~afterThree )
	\land ( takeFour = no \implies afterFour = afterThree ) \\
	( takeFive = yes \implies afterFive = (applyDelta~pendingDeltaFive)~afterFour )
	\land ( takeFive = no \implies afterFive = afterFour ) \\
	( takeSix = yes \implies crashNS = (applyDelta~pendingDeltaSix)~afterFive )
	\land ( takeSix = no \implies crashNS = afterFive ) \\
	\dom crashImages = refs~crashNS \\
	( \forall id : refs~crashNS @ \\
	\t1 ( id \in \dom syncedImages \land crashImages~id = syncedImages~id ) \lor \\
	\t1 ( id \notin \dom syncedImages \land id \in \dom dirtyChoices \land
	crashImages~id \in dirtyChoices~id ) ) \\
	visibleNS' = crashNS \land durableNS' = crashNS \land pendingCount' = pendingZero \\
	pendingDeltaOne' = pendingDeltaOne \land pendingDeltaTwo' = pendingDeltaTwo \land
	pendingDeltaThree' = pendingDeltaThree \land
	pendingDeltaFour' = pendingDeltaFour \land
	pendingDeltaFive' = pendingDeltaFive \land pendingDeltaSix' = pendingDeltaSix \land
	namespaceMutated' = namespaceMutated \\
	visibleImages' = crashImages \land syncedImages' = crashImages \land
	dirtyChoices' = \emptyset \\
	stageIDs' = \emptyset \land sealedIDs' = \emptyset \land closeOK' = \emptyset \\
	published' = published \cap \{ p : \dom crashNS; id : Identity |
	crashNS~p = file~id @ p \mapsto id \} \land
	publicationImages' = publicationImages \\
	owned' = \emptyset \land cleanupDue' = \emptyset \land
	cleanupPending' = \emptyset \land retainedResults' = \emptyset \land
	cleanupNotes' = \emptyset \\
	phase' = finished \land committed' = committed \land
	configSyncAck' = configSyncAck \land cancelled' = cancelled \land
	lockHeld' = no \land backupProgress' = backupProgress \land
	life' = ended \land reportState' = processDead \\
	revalidationPassed' = revalidationPassed \land
	confirmationAttempted' = confirmationAttempted \land
	cleanupWork' = \emptyset \land cleanupTried' = \emptyset \land
	cleanupSyncAttempted' = cleanupSyncAttempted \land releaseAttempted' = releaseAttempted
\end{schema}
% POINTS-CHECK END operation.powerloss

The literal authoritative Increment~1 operation inventory is:

\begin{zed}
	TxnStep \defs CandidateFailure \lor CreateStage \lor CreateStageFailure \lor \\
	\t1 WriteStage \lor WriteStageFailure \lor FileSyncStage \lor
	FileSyncStageFailure \lor CloseStage \lor CloseStageFailure \lor \\
	\t1 BeginRevalidation \lor RevalidationSuccess \lor RevalidationIOFailure \lor \\
	\t1 CommitAbsent \lor CommitAbsentFailure \lor ConfirmConfigSuccess \lor
	ConfirmConfigFailure \lor BeginCleanup \lor \\
	\t1 UnlinkOwned \lor CleanupAlreadyAbsent \lor CleanupUnlinkFailed \lor
	CleanupInspectionFailed \lor EndCleanupSweep \lor \\
	\t1 CleanupSyncSuccess \lor CleanupSyncFailure \lor ReleaseSuccess \lor
	ReleaseFailure \lor ReportDelivered \lor ReportFailed \lor \\
	\t1 FinishTransaction \lor ObserveCancellation \lor ProcessDeath \lor PowerLoss
\end{zed}

\subsection{Scenario monitor and generated wrappers}

The monitor is complete even though Scenario~A cannot exercise backup-order or
replacement injection.  The bad flags are independent latches.  No wrapper
assumes $MonitorSafety'$; a mutant must reach the corresponding bad state.

% POINTS-CHECK BEGIN monitor.types monitor
\begin{zed}
	SlotIndex == 1 \upto 6 \\
	DeathKind ::= noDeath | processDeathKind | powerLossKind \\
	OperationTag ::= opCandidateFailure | opCreateStage | opCreateStageFailure | \\
	\t1 opWriteStage | opWriteStageFailure | opFileSyncStage |
	opFileSyncStageFailure | opCloseStage | opCloseStageFailure | \\
	\t1 opBeginRevalidation | opRevalidationSuccess | opRevalidationIOFailure | \\
	\t1 opCommitAbsent | opCommitAbsentFailure | opConfirmConfigSuccess |
	opConfirmConfigFailure | opBeginCleanup | \\
	\t1 opUnlinkOwned | opCleanupAlreadyAbsent | opCleanupUnlinkFailed |
	opCleanupInspectionFailed | opEndCleanupSweep | \\
	\t1 opCleanupSyncSuccess | opCleanupSyncFailure | opReleaseSuccess |
	opReleaseFailure | opReportDelivered | opReportFailed | \\
	\t1 opFinishTransaction | opObserveCancellation | opProcessDeath | opPowerLoss
\end{zed}
% POINTS-CHECK END monitor.types

% POINTS-CHECK BEGIN monitor.vars monitor
\begin{schema}{CheckMonitor}
	seen : \power OperationTag \\
	linksInjected, accessInjected, replacementInjected : Bool \\
	cancellationAfterCommit : Bool \\
	deathKind : DeathKind \\
	lastCrashMask : \power SlotIndex \\
	preCrashActiveCount : PendingCount \\
	publicationSeen, publicationQualified : \power ( Path \cross Identity ) \\
	configPublicationSeen, revalidatedAtConfigPublication : Bool \\
	backupReadyAtConfigPublication : Bool \\
	replacementPath : \power Path \\
	replacementIdentity : \power Identity \\
	badRevalidation, badFilePublication, badBackupOrder : Bool \\
	badReplacementUnlink, badPrematureReport, badPostDeathStep : Bool \\
	badImageCoverage, badFinalImage, badRetainedCleanup : Bool
\end{schema}
% POINTS-CHECK END monitor.vars

% POINTS-CHECK BEGIN monitor.init monitor
\begin{schema}{MonitorInit}
	CheckMonitor'
	\where
	seen' = \emptyset \land linksInjected' = no \land accessInjected' = no \land
	replacementInjected' = no \land cancellationAfterCommit' = no \\
	deathKind' = noDeath \land lastCrashMask' = \emptyset \land
	preCrashActiveCount' = pendingZero \\
	publicationSeen' = \emptyset \land publicationQualified' = \emptyset \land
	configPublicationSeen' = no \land revalidatedAtConfigPublication' = no \land
	backupReadyAtConfigPublication' = no \\
	replacementPath' = \emptyset \land replacementIdentity' = \emptyset \\
	badRevalidation' = no \land badFilePublication' = no \land
	badBackupOrder' = no \land badReplacementUnlink' = no \land
	badPrematureReport' = no \land badPostDeathStep' = no \land
	badImageCoverage' = no \land badFinalImage' = no \land
	badRetainedCleanup' = no
\end{schema}
% POINTS-CHECK END monitor.init

The authoritative $MonitorInit$ schema is likewise source-only.  Its exact
predicate body is hashed and inlined into the sole generated initializer, so no
monitor initializer helper is exported as a ProZ component.

The generated wrapper supplies fixed $monitorOperation$, $monitorPath$, and
$monitorCrashMask$ locals, never ProB inputs.  The path is empty except for
cleanup operations and the crash mask is empty except for power loss.
$MonitorUpdate$ observes both transaction states but assigns only monitor state.

% POINTS-CHECK BEGIN monitor.update monitor
\begin{schema}{MonitorUpdate}
	\Delta TxnVars \\
	\Delta CheckMonitor \\
	monitorOperation : OperationTag \\
	monitorPath : \power Path \\
	monitorCrashMask : \power SlotIndex
	\where
	seen' = seen \cup \{ monitorOperation \} \\
	linksInjected' = linksInjected \land accessInjected' = accessInjected \land
	replacementInjected' = replacementInjected \\
	( cancellationAfterCommit' = yes \iff cancellationAfterCommit = yes \lor
	( monitorOperation = opObserveCancellation \land committed = yes ) ) \\
	( monitorOperation = opProcessDeath \implies deathKind' = processDeathKind \land
	preCrashActiveCount' = preCrashActiveCount \land lastCrashMask' = lastCrashMask ) \\
	( monitorOperation = opPowerLoss \implies deathKind' = powerLossKind \land
	preCrashActiveCount' = pendingCount \land lastCrashMask' = monitorCrashMask ) \\
	( monitorOperation \notin \{ opProcessDeath, opPowerLoss \} \implies
	deathKind' = deathKind \land preCrashActiveCount' = preCrashActiveCount \land
	lastCrashMask' = lastCrashMask ) \\
	publicationSeen' = publicationSeen \cup
	\{ p : Path; id : Identity |
	monitorOperation = opCommitAbsent \land p = configPath \land id = candidateID @
	p \mapsto id \} \\
	publicationQualified' = publicationQualified \cup
	\{ p : Path; id : Identity |
	monitorOperation = opCommitAbsent \land p = configPath \land id = candidateID \land
	id \in \dom visibleImages \cap \dom intendedImages \cap \dom syncedImages \land
	id \in sealedIDs \cap closeOK \land visibleImages~id = intendedImages~id \land
	syncedImages~id = intendedImages~id \land
	imageState~(intendedImages~id) = complete @ p \mapsto id \} \\
	( monitorOperation = opCommitAbsent \implies configPublicationSeen' = yes \land
	( revalidatedAtConfigPublication' = yes \iff
	revalidationPassed = yes \land phase = publishingConfig ) \land
	backupReadyAtConfigPublication' = yes ) \\
	( monitorOperation \neq opCommitAbsent \implies
	configPublicationSeen' = configPublicationSeen \land
	revalidatedAtConfigPublication' = revalidatedAtConfigPublication \land
	backupReadyAtConfigPublication' = backupReadyAtConfigPublication ) \\
	replacementPath' = replacementPath \land replacementIdentity' = replacementIdentity \\
	( badRevalidation' = yes \iff badRevalidation = yes \lor
	( monitorOperation = opCommitAbsent \land
	\lnot ( revalidationPassed = yes \land phase = publishingConfig ) ) ) \\
	( badFilePublication' = yes \iff badFilePublication = yes \lor \\
	\t1 ( monitorOperation = opCommitAbsent \land
	\lnot ( candidateID \in \dom visibleImages \cap \dom intendedImages
	\cap \dom syncedImages \land candidateID \in sealedIDs \cap closeOK \land
	visibleImages~candidateID = intendedImages~candidateID \land
	syncedImages~candidateID = intendedImages~candidateID \land
	imageState~(intendedImages~candidateID) = complete ) ) ) \\
	( badBackupOrder' = yes \iff badBackupOrder = yes \lor
	( monitorOperation = opCommitAbsent \land conflicts \neq \emptyset \land
	backupProgress \neq backupReady ) ) \\
	( badReplacementUnlink' = yes \iff badReplacementUnlink = yes \lor
	( monitorOperation = opUnlinkOwned \land monitorPath \neq \emptyset \land
	( \exists p : monitorPath @ p \in \dom visibleNS \land p \in \dom cleanupWork \land
	\lnot ( \{ p \mapsto file~(cleanupWork~p) \} \subseteq visibleNS ) ) ) ) \\
	( badPrematureReport' = yes \iff badPrematureReport = yes \lor \\
	\t1 ( monitorOperation \in \{ opReportDelivered, opReportFailed \} \land
	\lnot ( releaseAttempted = yes \land cleanupSyncAttempted = yes \land
		cleanupTried = \dom cleanupWork \land
			( committed = no \lor confirmationAttempted = yes ) ) ) ) \\
	( badPostDeathStep' = yes \iff badPostDeathStep = yes \lor life = ended ) \\
	( badImageCoverage' = yes \iff badImageCoverage = yes \lor
	( monitorOperation = opPowerLoss \land \dom visibleImages' \neq refs~visibleNS' ) ) \\
	( badFinalImage' = yes \iff badFinalImage = yes \lor \\
	\t1 ( monitorOperation = opPowerLoss \land
	( \exists p : Path; id : Identity | (p,id) \in publicationSeen \land
	\{ p \mapsto file~id \} \subseteq visibleNS' @ \\
	\t2 id \notin \dom visibleImages' \cap \dom publicationImages \lor \\
	\t2 ( id \in \dom visibleImages' \cap \dom publicationImages \land \\
	\t3 ( visibleImages'~id \neq publicationImages~id \lor
	imageState~(visibleImages'~id) \neq complete ) ) ) ) ) \\
	( badRetainedCleanup' = yes \iff badRetainedCleanup = yes \lor
	\dom retainedResults \cap \dom cleanupDue' \neq \emptyset \lor
	( monitorOperation = opUnlinkOwned \land
	monitorPath \cap \dom retainedResults \neq \emptyset ) )
\end{schema}
% POINTS-CHECK END monitor.update

% POINTS-CHECK BEGIN monitor.safety monitor
\begin{schema}{MonitorSafety}
	CheckMonitor
	\where
	badRevalidation = no \land badFilePublication = no \land
	badBackupOrder = no \land badReplacementUnlink = no \land
	badPrematureReport = no \land badPostDeathStep = no \land
	badImageCoverage = no \land badFinalImage = no \land
	badRetainedCleanup = no
\end{schema}
% POINTS-CHECK END monitor.safety

Scenario~A is the absent-file finite valuation.  The unused routing carriers
$ModelId$, $Provider$, and $Source$ are outside its transaction projection;
$JSONPointer$ remains because transaction state uses it for $conflicts$, and the
Scenario~A pointer carrier has cardinality one.  It has one candidate identity;
empty, prefix, and complete-candidate bytes; six extensional file images; and
exactly four deltas: idle, candidate-stage creation, config hard-link, and stage
removal.  Its exact finite carrier cardinalities are one JSON pointer, four
paths, one identity, one owner, one mode, three byte strings, four deltas, and
six file images.  Its initial visible/durable namespaces and all observation
image and metadata maps are empty; $expectedConfig=backupIDs=conflicts=\emptyset$;
$intendedImages$ maps the candidate to the complete candidate image; path roles
are the four distinct path atoms; authorization is accepted HTTPS without
insecure consent.  The generated operation set exports each Scenario~A operation
only as its corresponding \texttt{Check} wrapper, except that
$CleanupAlreadyAbsent$ remains a library relation but is not exported.  Scenario~A
has no environment transition that can make its sole cleanup path disappear
before the one-shot sweep, so exporting that wrapper would make exact nonzero
coverage impossible.  No backup, existing-file, metadata-injection,
external-disappearance, or replacement operation is introduced in Increment~1.

The required M1 anchor removes only the successful-revalidation call at
$CommitAbsent$.  The two M2 anchors remove only the file-sync calls at close and
publication qualification.  Their exact marker bodies include one tab, the
displayed predicate, two backslashes, and one LF.  M1 must reach
$badRevalidation=yes$; M2 must reach $badFilePublication=yes$ and, with
invariant stopping disabled, a power-loss state with $badFinalImage=yes$.

\section{On-disk configuration: what the tool writes, per source}

The tool writes, for each routed harness, only its \emph{route} (base URL +
auth, pointing at the chosen source's endpoint) and its \emph{default model} ---
plus, for Claude Code, the canonicalization override and the four tier knobs,
and, for pi/opencode, the native picker list. It never writes an agent
definition.

\textbf{Claude Code.} \texttt{.envrc.local} + \texttt{.claude/settings.local.json}.
Route: \texttt{ANTHROPIC\_BASE\_URL}, \texttt{ANTHROPIC\_AUTH\_TOKEN}. Default:
\texttt{ANTHROPIC\_MODEL} (verbatim, no $[1m]$ marker --- it breaks the
override-key match). Override: \texttt{modelOverrides}. Tier knobs:
\texttt{ANTHROPIC\_DEFAULT\_\{OPUS,SONNET,HAIKU,FABLE\}\_MODEL} and
\texttt{ANTHROPIC\_SMALL\_FAST\_MODEL}. No picker. A \texttt{.claude/agents/*.md}
with its own \texttt{model:} is a user agent definition --- out of scope.

\textbf{codex.} \texttt{config.toml}. Route: the matching
\texttt{[model\_providers.*]} table. Default: \texttt{model} and the single
\texttt{[agents] agent\_default\_subagent\_model}. No native picker list, no
per-agent files.

\textbf{pi.} \texttt{\~{}/.pi/agent/\{settings,models\}.json}. Route: the source
provider block + \texttt{defaultProvider}. Default: \texttt{defaultModel}.
Picker: the source provider's \texttt{models} array ($shown$) --- the
pattern-matched, routable subset. A \texttt{.pi/agents/*.md} frontmatter
\texttt{model:} is out of scope.

\textbf{opencode.} An explicit OpenCode V2 JSON file. Route: the
\texttt{providers.<name>} package, env, and settings leaves. Default: root
\texttt{model}. Registration: the \texttt{providers.<name>.models} map
($registered$), which is distinct from any optional picker view ($shown$).
Legacy singular \texttt{provider}, \texttt{options}, \texttt{small\_model}, and
\texttt{agent.<name>.model} are not written; agent declarations are protected.

The uniform rule: the tool writes the route, the default, and (for picker
harnesses) the picker view; it never writes an agent's model. A per-agent model
the user pinned still routes, because $routeVia$ is the one shared route for that
harness and the source serves whatever provider the pinned id names --- or it
does not, in which case that is the user's pin to fix, not the tool's.

\section{Sources as conforming instances}

Bifrost is one source: it exposes both protocols and both auth modes and serves
many providers by fanning out on the model-id prefix --- so a single Bifrost
source can carry all four harnesses, the single-source operator's case. A
direct source (direct OpenAI, direct Anthropic) typically exposes one protocol
and serves one provider; an aggregator (OpenRouter) may expose both protocols and
serve several providers without being a gateway. The multi-source operator
composes these: Claude Code routed through an Anthropic-capable source, the
OpenAI harnesses through whichever source they prefer. Nothing in the spec is
true only of a gateway; a re-implementation swaps or mixes sources without
changing the guarantees.

A source that can carry the whole fleet is one exposing both protocols; the spec
does not require any single such source to exist (the multi-source operator has
none), only that each routed harness's source speaks that harness's protocol.

\begin{axdef}
	fleetCapable : \power Source
	\where
	fleetCapable = \{ s : Source | \\
	\t1 anthropic \in srcProtocols~s \land openai \in srcProtocols~s \}
\end{axdef}

\section{Two real-world constraints: compatibility and credential terms}

\textbf{Compatibility (support, not permission).} Anthropic does not support
routing Claude Code to non-Claude models. This is a support and compatibility
statement, not a terms violation: with a non-Anthropic default on the Claude Code
route, tool use, extended thinking, and prompt caching are unvalidated and
degrade per model, so routability is necessary but not sufficient. Separately,
MCP tool search and fine-grained tool streaming silently turn off whenever Claude
Code is routed through any source other than Anthropic-direct; both are
recoverable via environment variables. These are side obligations of a routed
Claude Code, named here rather than encoded in the invariant.

\textbf{Credential terms (enforced).} The licensing line runs along the
credential, not the model: a Claude Pro/Max subscription is licensed for
Anthropic's own clients (the Claude apps and Claude Code), so spending it on a
third-party client is the violation --- for any model. Here $subauth$ is the
subscription-OAuth mode, and $Config$ plus $RouteOk$ permit it only where the
route is Anthropic-capable, which for a non-Claude harness would require that
harness to route through an Anthropic source it does not speak to --- excluded by
protocol coherence. $vkauth$ (per-host keys billing the upstream) is clean
everywhere.

\end{document}
